\documentclass[10pt]{article}
\usepackage[margin=1in]{geometry}
\usepackage[T1]{fontenc}
\usepackage{lmodern,amsmath,amssymb,amsthm,mathtools,microtype,array}
\usepackage[hidelinks]{hyperref}
\hypersetup{
 pdftitle={Exact Turan numbers of linear forests},
 pdfauthor={Qi Wu, Yong Lu},
 pdfsubject={Extremal graph theory},
 pdfkeywords={Turan number, linear forest, disjoint paths, Yuan-Zhang conjecture, Hamiltonian path}
}
\allowdisplaybreaks[2]
\setlength{\abovedisplayskip}{7pt plus 2pt minus 2pt}
\setlength{\belowdisplayskip}{7pt plus 2pt minus 2pt}
\setlength{\abovedisplayshortskip}{4pt plus 2pt minus 1pt}
\setlength{\belowdisplayshortskip}{4pt plus 2pt minus 1pt}
\newtheoremstyle{compactplain}{5pt}{5pt}{\itshape}{}%
 {\bfseries}{.}{.5em}{}
\theoremstyle{compactplain}
\newtheorem{theorem}{Theorem}[section]
\newtheorem{lemma}[theorem]{Lemma}
\newtheorem{proposition}[theorem]{Proposition}
\newtheorem{corollary}[theorem]{Corollary}
\newtheorem{conjecture}{Conjecture}[section]
\theoremstyle{remark}\newtheorem{remark}[theorem]{Remark}
\newcommand{\ex}{\operatorname{ex}}

\title{Exact Tur\'an numbers of linear forests\thanks{This work is supported by the National Natural Science Foundation of China (Nos.\ 12371348 and 12201258) and the High-Quality Science and Technology Cultivation Project of Jiangsu Normal University (No.\ JSNUGZL2026069).}}
\author{Qi Wu, Yong Lu\thanks{Corresponding author.}\\
School of Mathematics and Statistics, Jiangsu Normal University,\\
Xuzhou, Jiangsu 221116, People's Republic of China\\
Emails: \texttt{wuqimath@163.com, luyong@jsnu.edu.cn}}
\date{}

\begin{document}
\maketitle

\begin{abstract}
Let $H=\bigcup_{i=1}^s P_{k_i}$ be a linear forest with
$k_1\ge\cdots\ge k_s\ge2$, and set $m=\sum_{i=1}^s k_i$,
$a=\sum_{i=1}^s\lfloor k_i/2\rfloor$, and $S_j=\sum_{i=1}^j k_i$ for
$1\le j\le s$. Let $\varepsilon=1$ if every $k_i$ is odd, and
$\varepsilon=0$ otherwise. For integers $S\ge k\ge2$ and $n\ge0$, define
\[
 \Pi(n;S,k)=
 \begin{cases}
  \binom{n}{2}, & 0\le n<S-1,\\[2pt]
  \displaystyle\binom{S-1}{2}+u\binom{k-1}{2}+\binom{v}{2},
   & n=S-1+u(k-1)+v,\quad u\ge0,\quad 0\le v<k-1.
 \end{cases}
\]
Yuan and Zhang conjectured that whenever $k_s\ge3$, $k_1>3$, and
$n\ge m$,
\[
 \ex(n,H)=\max\left\{\Pi(n;S_j,k_j):1\le j\le s;\,
 (a-1)n-\binom{a}{2}+\varepsilon\right\}.
\]
In this paper we confirm this conjecture and show that the same formula
holds for every linear forest without isolated vertices that is not a
union of copies of $P_3$, thereby allowing components of order two as
well. Combined with the known exact formula for $tP_3$ and the trivial
identity $\ex(n,H)=\binom{n}{2}$ for $n<m$, this determines $\ex(n,H)$
for every linear forest without isolated vertices and all $n\ge0$.
\end{abstract}

\paragraph{Keywords:} Tur\'an number; linear forest; disjoint paths;
Yuan--Zhang conjecture; Hamiltonian path.

\paragraph{2020 Mathematics Subject Classification:} 05C35, 05C38.

\section{Introduction}\label{sec:intro}

All graphs considered in this paper are finite, simple, and undirected.
For a graph $H$, we say that a graph $G$ is \emph{$H$-free} if it
contains no subgraph isomorphic to $H$. The \emph{Tur\'an number}
$\ex(n,H)$ is the maximum number of edges in an $H$-free graph on $n$
vertices. We refer to the number of vertices of a graph $G$ as its
\emph{order}, denoted by $|G|$, and write $P_k$ for the path on $k$
vertices. Similarly, $K_q$ and $I_q$ denote the complete graph and the
edgeless graph on $q$ vertices, respectively, with the convention that
both are empty when $q=0$.

The Tur\'an problem for paths is a cornerstone of extremal graph theory.
Erd\H{o}s and Gallai~\cite{EG} established the classic linear upper bound
$\ex(n,P_k)\le\frac{k-2}{2}n$, and the exact value for all $n$ was later
determined independently by Faudree and Schelp~\cite{FS} and by
Kopylov~\cite{Kopylov}: writing $n=u(k-1)+v$ with $0\le v<k-1$, one has
$\ex(n,P_k)=u\binom{k-1}{2}+\binom{v}{2}$, attained by the disjoint
union of $u$ copies of $K_{k-1}$ and one copy of $K_v$.

A \emph{linear forest} is a graph whose connected components are paths.
For graphs $G$ and $H$, let $G\cup H$ denote their disjoint union, and
let $tG$ denote the disjoint union of $t$ copies of $G$ for any
nonnegative integer $t$. The Tur\'an problem for linear forests has
attracted considerable attention since the single-path case was settled.
Bushaw and Kettle~\cite{BushawKettle} determined $\ex(n,tP_k)$ for
sufficiently large $n$, and Lidick\'y, Liu, and Palmer~\cite{LLP}
extended their asymptotic result to arbitrary linear forests.

To describe the extremal constructions, recall that a \emph{matching} is
a set of pairwise vertex-disjoint edges, and $\nu(G)$ denotes the maximum
size of a matching in $G$. A set of pairwise nonadjacent vertices is
called \emph{independent}. The \emph{join} $G\vee H$ of two graphs is
obtained from $G\cup H$ by adding all edges between $V(G)$ and $V(H)$. We
call a path \emph{odd} or \emph{even} according to the parity of its
order. When $H$ is a linear forest without isolated vertices that is not
a union of copies of $P_3$, the large-$n$ extremal graph is
$K_{\nu(H)-1}\vee I_{n-\nu(H)+1}$, with one additional edge inside the
independent set when every component of $H$ is odd. The remaining
exceptional family $tP_3$ was solved for all $n$ independently by Campos
and Lopes~\cite{CL} and by Yuan and Zhang~\cite{YZtriples}, with the
latter authors also characterizing all extremal graphs.

For general linear forests at all values of $n$, Yuan and
Zhang~\cite{YZ} formulated an exact conjecture for
$H=\bigcup_{i=1}^s P_{k_i}$ with $k_1\ge\cdots\ge k_s\ge3$ and $k_1>3$.
They verified their formula for all $n$ whenever at most one component of
$H$ is odd, and characterized all extremal graphs when every component is
even. Subsequently, several specific families were settled for all $n$,
including $2P_5$~\cite{BielakKieliszek}, $2P_7$~\cite{LanQinShi},
$3P_5$~\cite{FengHu}, $2P_9$~\cite{ZhangWang}, and
$P_9\cup P_7$~\cite{FangYou}. Deng, Hou, and Zeng~\cite{DHZ} resolved
further three-path cases such as $3P_7$ and $2P_3\cup P_{2\ell+1}$, while
Lai, Ma, and Yan~\cite{LMY} settled linear forests whose components all
have order two or three. More recently, Dong, Ning, Yuan, and
Zhang~\cite{DNYZ} solved the remaining case of two odd paths, completing
the exact determination of $\ex(n,P_{k_1}\cup P_{k_2})$ for all
$k_1,k_2$. In a related direction, Chen, Yang, Yuan, and
Zhang~\cite{CYYZ} determined the maximum number of $r$-cliques in graphs
avoiding an even linear forest, and Zhang, Zhao, and Lu~\cite{ZZL}
studied Tur\'an problems forbidding both a fixed nonbipartite graph and a
linear forest.

Our proof also relies on degree-sum and dominating-path criteria for
spanning linear forests, or \emph{path factors}, developed by
Johansson~\cite{Johansson}, Chen, Tian, and Wei~\cite{CTW}, and Chiba and
Tsugaki~\cite{ChibaTsugaki}. What makes the all-$n$ problem challenging
is that different extremal constructions compete across different ranges
of $n$: disjoint unions of cliques dominate for small and intermediate
orders, whereas a join graph takes over once $n$ is large, with an extra
edge available inside the independent set in the all-odd case. Moreover,
when passing from connected graphs to general graphs, one must control
how components of $H$ can be distributed across different connected
components of $G$.

For integers $S\ge k\ge2$ and $n\ge0$, set
\begin{equation}\label{eq:prefix}
 \Pi(n;S,k)=
 \begin{cases}
  \binom{n}{2}, & 0\le n<S-1,\\[2pt]
  \displaystyle\binom{S-1}{2}+u\binom{k-1}{2}+\binom{v}{2},
   & n=S-1+u(k-1)+v,\quad u\ge0,\quad 0\le v<k-1.
 \end{cases}
\end{equation}
Note that the two expressions agree at $n=S-1$, and that
$\Pi(n;S,2)=\binom{S-1}{2}$ for all $n\ge S-1$. We also write
$J_a(n)=(a-1)n-\binom{a}{2}$ and
$B_m(n)=\binom{m-1}{2}+\lfloor(n-m+1)/2\rfloor$.

Given $H=\bigcup_{i=1}^s P_{k_i}$ with $k_1\ge\cdots\ge k_s\ge2$, put
$m=\sum_{i=1}^s k_i$, $a=\sum_{i=1}^s\lfloor k_i/2\rfloor=\nu(H)$, and
$S_j=\sum_{i=1}^j k_i$ for $1\le j\le s$. For each $j$, the quantity
$\Pi(n;S_j,k_j)$ is the edge count of a clique of order $S_j-1$ together
with a $P_{k_j}$-free disjoint union of cliques, while $J_a(n)$ is the
edge count of $K_{a-1}\vee I_{n-a+1}$ (with one extra edge in the
independent set when every $k_i$ is odd).

\begin{conjecture}[Yuan and Zhang~\cite{YZ}]
Let $H=\bigcup_{i=1}^s P_{k_i}$ with $k_1\ge\cdots\ge k_s\ge3$ and
$k_1>3$. Put $m=\sum_{i=1}^s k_i$,
$a=\sum_{i=1}^s\lfloor k_i/2\rfloor$, and $S_j=\sum_{i=1}^j k_i$. Let
$\varepsilon=1$ if every $k_i$ is odd, and $\varepsilon=0$ otherwise.
Then, for every $n\ge m$,
\begin{equation}\label{eq:formula}
 \ex(n,H)=\max\bigl\{\Pi(n;S_j,k_j):1\le j\le s;\,
 J_a(n)+\varepsilon\bigr\}.
\end{equation}
\end{conjecture}

Our main theorem settles this conjecture in full and extends
\eqref{eq:formula} to allow $P_2$ components. For completeness, we
include the known matching and $tP_3$ cases in the statement.

\begin{theorem}\label{thm:main}
Let $H=\bigcup_{i=1}^s P_{k_i}$ with $k_1\ge\cdots\ge k_s\ge2$. Put
$m=\sum_{i=1}^s k_i$, $a=\sum_{i=1}^s\lfloor k_i/2\rfloor$, and
$S_j=\sum_{i=1}^j k_i$. Let $\varepsilon=1$ if every $k_i$ is odd, and
$\varepsilon=0$ otherwise. If $H$ is not a union of copies of $P_3$, then
\eqref{eq:formula} holds for every $n\ge m$. If $H=tP_3$ with $t\ge1$,
then for every $n\ge3t$,
\begin{equation}\label{eq:triples}
 \ex(n,tP_3)=\max\left\{
 \binom{3t-1}{2}+\left\lfloor\frac{n-3t+1}{2}\right\rfloor,\,
 (t-1)n-\binom{t}{2}+\left\lfloor\frac{n-t+1}{2}\right\rfloor
 \right\}.
\end{equation}
Finally, $\ex(n,H)=\binom{n}{2}$ whenever $0\le n<m$.
\end{theorem}

Formula~\eqref{eq:triples} was established independently by Campos and
Lopes~\cite{CL} and by Yuan and Zhang~\cite[Theorem~1.6]{YZtriples}; the
difference between its second and first terms is $(t-1)(n-5t+1)$, so the
maximum agrees with the piecewise expression in~\cite{YZtriples}. When
$H=P_2$, we have $\ex(n,P_2)=0$ for $n\ge2$, and general matchings
$tP_2$ are covered by the Erd\H{o}s--Gallai theorem~\cite{EG}. We remark
that Yuan and Zhang~\cite{YZ} also asked for a full characterization of
the extremal graphs; here we determine the exact Tur\'an numbers and
verify the lower-bound constructions below.

\paragraph{Lower bounds.}
Fix $j\in\{1,\ldots,s\}$ and $n\ge m$, and consider the graph
$K_{S_j-1}\cup uK_{k_j-1}\cup K_v$, where $u$ and $v$ are given by
\eqref{eq:prefix}. Every component among the first $j$ paths of $H$ has
order at least $k_j$, so any copy of $H$ would have to place all $j$ of
them inside the clique $K_{S_j-1}$, which is impossible since they need
$S_j$ vertices.

Next, recall that a \emph{vertex cover} is a vertex set meeting every
edge of a graph. The join $K_{a-1}\vee I_{n-a+1}$ has $J_a(n)$ edges and
a vertex cover of size $a-1$, whereas every vertex cover of $H$ has size
at least $\nu(H)=a$; so this join is $H$-free. Suppose now that every
component of $H$ is odd, and add a single edge inside $I_{n-a+1}$. Any
odd path $P_{2b+1}$ using $c$ vertices of $K_{a-1}$ breaks into at most
$c+1$ segments outside $K_{a-1}$, at most one of which can have two
vertices (using the added edge). Its order is therefore at most $2c+2$,
which forces $c\ge b$. Summing over all components of $H$ shows that any
copy of $H$ would require at least $a$ vertices of $K_{a-1}$, a
contradiction.

For $tP_3$, the two extremal graphs $K_{3t-1}\cup M_{n-3t+1}$ and
$K_{t-1}\vee M_{n-t+1}$ (where $M_q$ consists of a matching of size
$\lfloor q/2\rfloor$ plus an isolated vertex when $q$ is odd) are easily
seen to be $tP_3$-free.

Isolated vertices can be accommodated immediately.

\begin{corollary}
Let $F=H\cup zP_1$, where $z\ge0$ and $H$ is a linear forest with at
least one edge and no isolated vertices. Then
\[
 \ex(n,F)=
 \begin{cases}
  \binom{n}{2}, & n<|F|,\\
  \ex(n,H), & n\ge|F|,
 \end{cases}
\]
where $\ex(n,H)$ is given by Theorem~\ref{thm:main}.
\end{corollary}
\begin{proof}
For $n<|F|$ the complete graph $K_n$ is $F$-free. If $n\ge|H|+z$, then
any copy of $H$ leaves at least $z$ unused vertices that can serve as the
isolated components of $F$, so a graph on $n$ vertices is $F$-free if
and only if it is $H$-free.
\end{proof}

If a linear forest consists solely of isolated vertices, then no graph of
order at least $|F|$ avoids it. From now on we assume that every
forbidden linear forest has at least one edge.

The rest of the paper is organized as follows. Section~\ref{sec:prelim}
gathers the preliminary tools (the connected path bound of
Balister, Gy\H{o}ri, Lehel, and Schelp, path-factor criteria, and
Hamiltonian-path closure) and resolves the connected case when $H$ has at
most two odd components. Section~\ref{sec:coretools} develops the
$a$-core decomposition, an ordered odd-path deletion lemma, and edge
estimates for the core and residual components. In
Section~\ref{sec:connected} we combine a P\'osa-type structural analysis
with Hall's theorem to complete the proof for connected graphs. Finally,
Section~\ref{sec:main} extends the connected bound to arbitrary graphs
via a component-allocation inequality, completing the proof of
Theorem~\ref{thm:main}.

\section{Preliminary lemmas and the first connected cases}\label{sec:prelim}

For a graph $G$, we write $V(G)$ and $E(G)$ for its vertex and edge
sets, and $e(G)=|E(G)|$. Set $L_q(n)=qn-\binom{q+1}{2}$, so that
$J_a(n)=L_{a-1}(n)$. For $v\in V(G)$, let $N_G(v)$ denote the
\emph{neighborhood} of $v$ and put $d_G(v)=|N_G(v)|$. When
$X\subseteq V(G)$ (or $X$ is an induced subgraph of $G$), we write
$N_X(v)=N_G(v)\cap X$ and $d_X(v)=|N_X(v)|$, dropping the subscript when
the ambient graph is clear. We denote by $G[X]$ the subgraph induced by
$X$ and set $G-X=G[V(G)\setminus X]$; for a subgraph $Q$ we write
$G-Q=G-V(Q)$, and for a single vertex $v$ we abbreviate $G-\{v\}$ to
$G-v$. We use $\dot\cup$ to emphasize that a union of sets or graphs is
disjoint.

By the theorem of Faudree and Schelp~\cite{FS} (see also
\cite[Theorem~1.2]{BGLS}), if $n=u(k-1)+v$ with $0\le v<k-1$ and
$k\ge3$, then $\ex(n,P_k)=u\binom{k-1}{2}+\binom{v}{2}$; in particular,
$\ex(n,P_k)=\binom{n}{2}$ for $n<k$. Balister, Gy\H{o}ri, Lehel, and
Schelp~\cite[Theorem~1.3]{BGLS} proved the following connected
counterpart: if $X$ is a connected $P_{k+1}$-free graph of order
$x>k\ge3$, then
\begin{equation}\label{eq:bgls}
 e(X)\le\max\left\{
 \binom{k-1}{2}+x-k+1,\quad
 \binom{k-b}{2}+b(x-k+b)\right\},
 \qquad b=\left\lfloor\frac{k-1}{2}\right\rfloor.
\end{equation}
Moreover, equality holds only when
$X\cong K_s\vee\bigl(K_{k-2s}\mathbin{\dot\cup}I_{x-k+s}\bigr)$ for some
$s\in\{1,b\}$ (these two graphs and bounds coincide for $k\in\{3,4\}$).

For a noncomplete graph $X$, define
$\sigma_2(X)=\min\{d_X(u)+d_X(v):uv\notin E(X)\}$, and set
$\sigma_2(X)=\infty$ if $X$ is complete. We will use the following
consequence of the path-factor theorem of Chen, Tian, and Wei~\cite{CTW}
(see also Magnant and Ozeki~\cite[Theorem~1.2]{MO}).

\begin{theorem}\label{thm:ctw}
Let $t\ge2$, let $a_1,\ldots,a_t\ge2$, and put $n=\sum_{i=1}^t a_i$. Let
$\lambda_{\rm e}$ be the number of even integers among
$a_1,\ldots,a_t$. If a graph $X$ of order $n$ satisfies
$\sigma_2(X)\ge n-t+\lambda_{\rm e}-2$, then $X$ contains a spanning
path factor with component orders $a_1,\ldots,a_t$.
\end{theorem}

We repeatedly use a simple convexity observation: if
$E(x)=\max\{C,qx-B\}$ with $q>0$ and $x_0=(C+B)/q$, then any linear
function whose slope lies in $[0,q]$ and whose value at $x_0$ is less
than $C$ stays strictly below $E(x)$ for all $x$ (compare with $C$ for
$x\le x_0$ and with $qx-B$ for $x\ge x_0$).

A path $P$ in a graph is \emph{$d$-dominating} if every vertex outside
$P$ has at least $d$ neighbors on $P$. Chiba and
Tsugaki~\cite[Lemma~2.2]{ChibaTsugaki} recorded the dominating-path
version of Johansson's path-factor theorem~\cite{Johansson}. The next
lemma lowers that degree threshold by one in the specific configuration
needed here.

\begin{lemma}\label{lem:path-completion}
Let $F$ be a linear forest with exactly two odd components (each of
order at least three) and at least one even component, and put
$|F|=2a+2$. If $P=v_1\cdots v_{2a+1}$ is a path and $z\notin V(P)$
satisfies $d_P(z)\ge a-1$, then $P\cup\{z\}$ contains a spanning copy of
$F$.
\end{lemma}
\begin{proof}
Merging all even components of $F$ into a single path of order $2c$, we
write the two odd component orders as $2b+1$ and $2d+1$; it suffices to
find a spanning copy of $P_{2c}\cup P_{2b+1}\cup P_{2d+1}$, since
$P_{2c}$ can then be split into the required even components. Here
$a=b+c+d$ with $b,c,d\ge1$. Because exactly one of the three target
paths is even, Johansson's dominating-path threshold~\cite[Lemma~2.2]{ChibaTsugaki}
is $a$, which handles $d_P(z)\ge a$. We may therefore assume that
$d_P(z)=a-1$.

Write $N=\{i:zv_i\in E\}$, so $|N|=a-1$, and suppose for a contradiction
that no such spanning factor exists. Fix an ordering
$\ell_1,\ell_2,\ell_3$ of the three path orders, and set $S_0=0$ and
$S_i=\sum_{j\le i}\ell_j$. Clearly $N$ cannot contain two consecutive
indices, as $z$ could otherwise be inserted into $P$ to form a spanning
path. Likewise, if $S_{i-1}+1\in N$ or $S_i-1\in N$, then $z$ together
with $v_{S_{i-1}+1}\cdots v_{S_i-1}$ forms a path of order $\ell_i$
while the remaining two intervals supply the other two components. Thus
$N$ has no consecutive elements and avoids $S_{i-1}+1$ and $S_i-1$ for
every $i\in\{1,2,3\}$.

The only available indices are $S_1,S_2$ and the three intervals
$I_i=\{S_{i-1}+2,\ldots,S_i-2\}$. For the even path, the corresponding
interval is empty when $c=1$, and when $c\ge2$ it has $2c-3$ positions,
allowing at most $c-1$ nonconsecutive elements (with a unique extremal
set); in either case it contributes at most $c-1$ to $|N|$. Similarly,
the two odd intervals contribute at most $b-1$ and $d-1$, respectively.
This yields $|N|\le2+(c-1)+(b-1)+(d-1)=a-1$. Equality must therefore
hold, meaning that $S_1,S_2\in N$ and every interval $I_i$ attains its
maximum.

Placing the even path first forces $2,4,\ldots,2c\in N$ (which is simply
$S_1=2$ when $c=1$), while placing it last forces
$2a-2c+2,2a-2c+4,\ldots,2a\in N$. In particular, $\min N=2$ and
$\max N=2a$. On the other hand, placing $P_{2b+1}$ first excludes both
$2b$ and $2b+2$ from $N$.

Listing $N=\{s_1<\cdots<s_{a-1}\}$, we have $s_{i+1}-s_i\ge2$ for all
$i$ and $\sum_{i=1}^{a-2}(s_{i+1}-s_i-2)=2$, so either one gap equals
$4$ or two gaps equal $3$. A single gap of $4$ would force $N$ to
contain every even index in $[2,2a]$ except one, contradicting that both
$2b$ and $2b+2$ are absent. Thus there are two gaps of size $3$. Let the
first run from $2h$ to $2h+3$. Since $2,4,\ldots,2c\in N$, we have
$h\ge c$, and because all even positions up to $2h$ lie in $N$, we have
$x=2h-2c+2\in N$ as well as $y=2h+3\in N$. Now $v_{x+1}\cdots v_{y-1}$
is a path of order $2c$, and the remaining vertices form the path
$v_1\cdots v_x z v_y\cdots v_{2a+1}$ of order $2b+2d+2$, which splits
into $P_{2b+1}$ and $P_{2d+1}$, a contradiction.
\end{proof}

Next, we collect Lemmas~1.4, 2.3, and 2.4 of Balister et
al.~\cite{BGLS} in the notation needed below; the lower bound $p\ge1$
when $s=0$ follows directly from connectedness.

\begin{lemma}\label{lem:bgls-weights}
Let $G$ be a connected graph, let $P$ be a longest path in $G$ of order
$k$, and set $Y=V(G)\setminus V(P)\ne\varnothing$. For each $y\in Y$,
let $s_y=d_P(y)$ and let $p_y$ be the number of edges in a longest path
of $G[Y]$ starting at $y$. Putting $w_y=s_y+p_y/2$ and
$\lambda=\binom{k}{2}-e(G[V(P)])$, we have
\begin{equation}\label{eq:weights}
 e(G)\le e(G[V(P)])+\sum_{y\in Y}w_y.
\end{equation}
Moreover, writing $s=s_y$ and $p=p_y$,
\[
\begin{array}{ll}
 s>0: & 2s+2p\le k-1,\quad
 \lambda\ge(p+1)(k-p-2)+\binom{s-1}{2},\\[1mm]
 s=0: & 1\le p\le k-3,\quad
 \lambda\ge\dfrac{k-1}{2}\bigl(\lfloor p/2\rfloor+1\bigr).
\end{array}
\]
\end{lemma}

For integers $1\le t<\ell/2$, define
$F_\ell(t)=\binom{\ell-t}{2}+t(t+1)$ and
$B(\ell,t)=\max\{F_\ell(t),F_\ell(\lfloor(\ell-1)/2\rfloor)\}$; by
convexity, $B(\ell,t)=\max_{t\le j<\ell/2}F_\ell(j)$. A
\emph{Hamiltonian path} and a \emph{Hamiltonian cycle} are a spanning
path and a spanning cycle, respectively. By Balister et
al.~\cite[Lemma~2.2, equations~(2.1)--(2.2)]{BGLS}, if a connected graph
$R$ on $k+1$ vertices has no Hamiltonian path, $R-v$ has a Hamiltonian
path, and $d_R(v)=d$, then $e(R)\le B(k,d)$. Choosing a removable vertex
of maximum degree sharpens this bound as follows.

\begin{lemma}\label{lem:omitted-vertex}
Let $R$ be a connected graph on $k+1$ vertices that contains $P_k$ but
has no Hamiltonian path. Call a vertex $v\in V(R)$ \emph{removable} if
$R-v$ has a Hamiltonian path, and let $d$ be the maximum degree of a
removable vertex. If $1\le d<\lfloor(k-1)/2\rfloor$, then
\begin{equation}\label{eq:omitted-bound}
 e(R)\le\max\bigl\{F_k(d),\,B(k-1,d)+\lfloor k/2\rfloor\bigr\}.
\end{equation}
If $d=\lfloor(k-1)/2\rfloor$, then $e(R)\le F_k(d)$.
\end{lemma}
\begin{proof}
Pick a removable vertex $v$ with $d_R(v)=d$, and let
$P=x_1\cdots x_k$ be a Hamiltonian path of $R-v$. Write
$N_R(v)=\{x_{a_1},\ldots,x_{a_d}\}$ with $a_1<\cdots<a_d$. Since $R$ has
no Hamiltonian path, $v$ is nonadjacent to $x_1$ and $x_k$ and has no
two consecutive neighbors on $P$. Set
$T=\{v,x_1,x_{a_1+1},\ldots,x_{a_d+1}\}$. An edge $x_1x_{a_i+1}$ would
give the Hamiltonian path
$vx_{a_i}x_{a_i-1}\cdots x_1x_{a_i+1}\cdots x_k$, while an edge
$x_{a_i+1}x_{a_j+1}$ with $i<j$ would give
$x_1\cdots x_{a_i}vx_{a_j}x_{a_j-1}\cdots x_{a_i+1}x_{a_j+1}\cdots x_k$;
so $T$ is independent.

In fact, adding an edge between any two vertices of $T\setminus\{v\}$
creates a Hamiltonian path by the same rerouting. By the Bondy--Chv\'atal
closure theorem for Hamiltonian paths (recalled in the next subsection),
every pair in $T\setminus\{v\}$ has degree sum at most $k-1$, so at most
one vertex of $T\setminus\{v\}$ has degree exceeding $\lfloor k/2\rfloor$.

Suppose first that some $w\in T\setminus\{v\}$ satisfies
$d<d_R(w)\le\lfloor k/2\rfloor$. Form a graph $R'$ on
$V(R)\setminus\{w\}$ by deleting $w$ from $R$ and, if $w$ is an internal
vertex of $P$, adding the edge between its two neighbors on $P$ (if not
already present) and marking that added edge. Then $R'-v$ still has a
Hamiltonian path, $d_{R'}(v)=d$ (since $vw\notin E(R)$), and $R'$ is
connected. Moreover, $R'$ has no Hamiltonian path: a Hamiltonian path
using the marked edge could be expanded through $w$ to a Hamiltonian path
of $R$, whereas one avoiding it (or occurring when $w$ is an endpoint of
$P$) would be a Hamiltonian path of $R-w$, making $w$ removable despite
$d_R(w)>d$. Applying the BGLS estimate to $R'$ gives
$e(R)=e(R-w)+d_R(w)\le e(R')+d_R(w)\le B(k-1,d)+\lfloor k/2\rfloor$.

If no such $w$ exists, then all vertices of $T\setminus\{v\}$ except at
most one have degree at most $d$, so $T$ contains at least $d+1$ vertices
of degree at most $d$. Counting edges outside those $d+1$ vertices and
edges incident with them yields
$e(R)\le\binom{k-d}{2}+d(d+1)=F_k(d)$, establishing
\eqref{eq:omitted-bound}.

Finally, suppose $d=\lfloor(k-1)/2\rfloor$. When $k=2q+1$, we have
$d=q=\lfloor k/2\rfloor$, so the first case cannot arise and the second
gives $e(R)\le F_{2q+1}(q)$. When $k=2q$, either some vertex of
$T\setminus\{v\}$ has degree $q$, in which case the construction of $R'$
yields $e(R)\le B(2q-1,q-1)+q=F_{2q}(q-1)$, or no such vertex exists and
the low-degree bound gives $F_{2q}(q-1)$ directly.
\end{proof}

\begin{lemma}\label{lem:p2-attachment}
Let $H$ be a nonempty linear forest with all components of order at
least two, and put $Q=\nu(H)$, $h=|H|$, and $o$ for the number of odd
components of $H$. If $G$ is an $(H\cup P_2)$-free graph containing a
path $P=v_1\cdots v_{h+1}$, then $V(G)\setminus V(P)$ is independent and
every $x\in V(G)\setminus V(P)$ satisfies $d_P(x)\le Q-o$.
\end{lemma}
\begin{proof}
Any edge outside $P$ together with a copy of $H$ inside $P$ would yield
$H\cup P_2$, so $V(G)\setminus V(P)$ is independent.

Fix $x\notin V(P)$. Two neighbors of $x$ on $P$ cannot be consecutive
(otherwise $x$ inserts into $P$ to give $P_{h+2}$, which contains
$H\cup P_2$), nor can their indices differ by $3$ (otherwise we take the
edge between the two intervening vertices as $P_2$ and route the rest of
$P$ through $x$, leaving a path of order $h$ that contains $H$).

Order the components of $H$ arbitrarily, with partial order sums
$0=S_0<S_1<\cdots<S_s=h$. For every $0\le j\le s$, we must have
$xv_{S_j+1}\notin E(G)$, since otherwise $xv_{S_j+1}$ serves as $P_2$
while the first $j$ components embed into $v_1\cdots v_{S_j}$ and the
remaining ones into $v_{S_j+2}\cdots v_{h+1}$.

Now examine a component of order $k$ starting after index $S$; all
possible neighbors of $x$ in its block lie in $[S+2,S+k]$. If $k=2b$,
having no consecutive neighbors allows at most $b$ neighbors in this
interval. Suppose $k=2b+1$. Then $S+3$ and $S+k-1$ are also forbidden:
if $xv_{S+3}\in E(G)$, we take $v_{S+1}v_{S+2}$ as $P_2$,
$xv_{S+3}\cdots v_{S+k+1}$ as the current component, and fit the earlier
and later components on $v_1\cdots v_S$ and
$v_{S+k+2}\cdots v_{h+1}$; symmetrically, $xv_{S+k-1}\in E(G)$ allows us
to use $v_{S+k}v_{S+k+1}$ as $P_2$ and $v_{S+1}\cdots v_{S+k-1}x$ as the
current component. If $x$ had $b$ neighbors in $[S+2,S+k]$, then listing
their indices in increasing order would give gaps of size at least $2$
and never $3$, within a span of at most $2b-1$. A span of $2b-1$ would
force one gap of size $3$, so the span must be $2b-2$ and every gap must
equal $2$. The $b$ neighbors would then all have the same parity, which
is impossible because the two excluded indices $S+3$ and $S+2b$ have
opposite parity. Thus each odd component contributes at most $b-1$
neighbors, and summing over all components gives $d_P(x)\le Q-o$.
\end{proof}

\begin{lemma}\label{lem:low-attachment}
Let $q\ge3$, and let $G$ be a connected graph of order $n\ge2q+4$ whose
longest path has order $k=2q+3$. Choose a longest path $P$ maximizing
$d=\max_{y\notin V(P)}d_P(y)$. If $d\le q-1$, then
$e(G)\le\max\{\binom{k}{2}+\lfloor(n-k)/2\rfloor,L_q(n)\}$.
\end{lemma}
\begin{proof}
Set $Y=V(G)\setminus V(P)$, $N=|Y|$, $K=\binom{k}{2}$, and
$\lambda=K-e(G[P])$. Also put $C=K-L_q(k)=(q+2)(q+3)/2$ and
$z=C/(q-1/2)=(q+2)(q+3)/(2q-1)$.

For each $y\in Y$, let $s,p$, and $w=s+p/2$ be as in
Lemma~\ref{lem:bgls-weights}. Then $1/2\le w\le q$ and
\[
 \lambda\ge f(s,p):=
 \begin{cases}
  (p+1)(2q+1-p)+\binom{s-1}{2}, & s>0,\\
  (q+1)(\lfloor p/2\rfloor+1), & s=0.
 \end{cases}
\]
If $\lambda\ge C$, then \eqref{eq:weights} immediately gives
$e(G)\le L_q(n)$, so assume $\lambda<C$. For $1\le s\le q-1$, we cannot
have $p=q+1-s$ because $f(s,q+1-s)-C=(q^2-q-s^2+s-4)/2\ge q-3\ge0$;
hence $p\le q-s$.

Pick $x\in Y$ with $d_P(x)=d$. By the choice of $P$, the maximum degree
of a removable vertex in $G[P\cup\{x\}]$ is $d$. Applying
Lemma~\ref{lem:omitted-vertex} and noting that
$F_k(d)-F_{k-1}(d)-(q+1)=q+1-d>0$, we get
\[
 \lambda\ge\min\{\alpha(d),\beta(d)\},\qquad
 \alpha(s)=\frac{s(4q+5-3s)}{2},\quad
 \beta(s)=\frac{q^2+3q+2}{2}+s.
\]
For every $1\le s\le d\le q-1$, we have
$\min\{\alpha(d),\beta(d)\}\ge\min\{\alpha(s),\beta(s)\}$: for $q\ge4$
this follows from the concavity of $\alpha$, the monotonicity of
$\beta$, and $\alpha(q-1)-\beta(q-1)=q-4\ge0$, while for $q=3$ direct
evaluation at $s=1,2$ gives $\alpha=7,11$ and $\beta=11,12$.

We claim that $\lambda\ge z(w-1/2)$ for every $y\in Y$. Consider first
$s=0$. If $p=2u$ with $1\le u\le q$, then $w=u$ and
$\lambda-z(w-1/2)\ge D_{\rm e}(u):=(q+1)(u+1)-z(u-1/2)$; if $p=2u+1$
with $0\le u\le q-1$, then $w=u+1/2$ and
$\lambda-z(w-1/2)\ge D_{\rm o}(u):=(q+1)(u+1)-zu$. Both linear functions
have slope $q+1-z$, which is positive for $q\ge6$, where
$D_{\rm e}(1)=(7q^2-q-10)/(2(2q-1))>0$ and $D_{\rm o}(0)=q+1$. For
$q=3,4,5$, checking the right endpoints gives the pairs $(1,0)$,
$(4,2)$, and $(8,46/9)$, proving the claim for $s=0$.

Now suppose $s>0$. The quantity $\beta(s)-z(q+s-1)/2$ decreases in $s$
and equals $(q-4)(q-3)/(2(2q-1))\ge0$ at $s=q-1$, so
$\beta(s)\ge z(w-1/2)$ since $w\le(q+s)/2$. Furthermore,
$\lambda\ge\min\{\beta,(3\alpha+f)/4\}$. Define
$\Psi(s,p)=(3\alpha(s)+f(s,p))/4-z(s+p/2-1/2)$, which is concave in
$(s,p)$. Every feasible integer pair $(s,p)$ except $(q-1,1)$ lies in
the convex hull of $(1,0)$, $(q-1,0)$, $(q-2,2)$, and $(1,q-1)$. At
these four corners, $4(2q-1)\Psi$ equals $2(7q^2-5q-8)$,
$4(q^2-5q+11)$, $2(q-3)(8q-11)$, and $(q-3)(5q+1)$, respectively, all
nonnegative. At $(q-1,1)$, we have
$\alpha(q-1)-z(q-1)=5(q-4)(q-1)/(2(2q-1))\ge0$ for $q\ge4$, and
$f(2,1)=12=z(2+1/2-1/2)$ for $q=3$. Thus $\lambda\ge z(w-1/2)$ in all
cases.

Writing $\bar w=N^{-1}\sum_{y\in Y}w_y$, inequality \eqref{eq:weights}
gives $e(G)\le K-\lambda+N\bar w\le K+z/2+(N-z)\bar w$. When $N\le z$,
using $\bar w\ge1/2$ yields $e(G)\le K+N/2$, and hence
$e(G)\le K+\lfloor N/2\rfloor$ by integrality; when $N\ge z$, using
$\bar w\le q$ yields $e(G)\le L_q(n)$.
\end{proof}

\subsection{Hall's theorem and the closure theorem}

We recall two standard tools. First, by Hall's theorem (see
Diestel~\cite[Chapter~2]{Diestel}), a bipartite graph with bipartition
$(X,Y)$ has a matching covering $X$ if and only if $|N(S)|\ge|S|$ for
every $S\subseteq X$, where $N(S)$ is the set of neighbors of $S$ in
$Y$.

Second, by the Bondy--Chv\'atal closure theorem~\cite{BondyChvatal}, if
$xy$ is a nonedge of an $m$-vertex graph, then adding $xy$ preserves the
existence of a Hamiltonian cycle when $d(x)+d(y)\ge m$, a Hamiltonian
path when $d(x)+d(y)\ge m-1$, and a Hamiltonian path with two prescribed
endpoints when $d(x)+d(y)\ge m+1$ (the path statements follow from the
cycle version by adjoining an auxiliary vertex and deleting it
afterward). A vertex adjacent to all other vertices of a graph is called
\emph{universal}.

\subsection{Connected graphs with at most two odd components}

We now treat connected graphs when the forbidden linear forest has at
most two odd components. The first lemma will also be needed in
Section~\ref{sec:connected}.

\begin{lemma}\label{lem:two-odd-defect}
Let $F=P_{2q+1}\cup P_{2s+1}$ with $q\ge s\ge1$, $q+s=a\ge3$, and
$(q,s)\ne(1,1)$. Let $Q$ be a graph with vertex partition
$V(Q)=S\dot\cup T$, where $|S|=a+2$, $|T|=a$, and the vertices of $T$
are linearly ordered. For each $t\in T$, let $d_T^+(t)$ denote the
number of neighbors of $t$ appearing later in this order. If nonnegative
integers $\delta_t$ satisfy $d_S(t)+d_T^+(t)\ge a-1-\delta_t$ for all
$t\in T$ and $\lambda+\Delta\le1$, where
$\lambda=\binom{a+2}{2}-e(Q[S])$ and $\Delta=\sum_{t\in T}\delta_t$,
then $Q$ contains a spanning copy of $F$.
\end{lemma}
\begin{proof}
If $\Delta=0$, then every vertex of $T$ has degree at least $a-1$, while
$Q[S]$ misses at most one edge and thus has minimum degree at least $a$;
every nonadjacent pair in $Q$ therefore has degree sum at least $2a-2$,
and Theorem~\ref{thm:ctw} yields a spanning copy of $F$. We may thus
assume $\lambda=0$ and $\Delta=1$, so $Q[S]$ is complete and at most one
vertex of $T$ has degree defect $1$.

Suppose that $Q$ has no Hamiltonian path, and apply Hamiltonian-path
closure iteratively. Whenever a vertex $x\in S$ has a neighbor in $T$,
we have $d_Q(x)\ge a+2$, so closure first joins $x$ to all vertices of
$T$ of degree at least $a-1$ (raising $d_Q(x)$ to at least $2a$) and
then also to the possible vertex of degree $a-2$; thus $x$ becomes
universal. Let $A\subseteq S$ be the set of vertices of $S$ that become
universal under this procedure, and put $B=S\setminus A$ and $p=|B|$. By
construction, no edge runs between $B$ and $T$.

Let $u$ and $v$ be the last two vertices of $T$. Their forward degrees
sum to at most $1$, so adding their degree inequalities gives
$2(a-1)-\delta_u-\delta_v\le2|A|+1$. Using $\delta_u+\delta_v\le1$ and
$|A|=a+2-p$ yields $p\le4$. On the other hand, if $p\le1$, then the
vertices of $B\cup T$ as singletons form at most $|A|+1$ path components
that can be linked via the universal set $A$ into a Hamiltonian path.
Hence $2\le p\le4$.

If $p=2$, the edge inside $Q[B]=K_2$ already reduces $B\cup T$ to a
spanning linear forest with $|A|+1$ components, again giving a
Hamiltonian path. If $p=4$, summing the degree inequalities over $T$
gives $e(Q[T])\ge a-1\ge2$, and a spanning path of $Q[B]=K_4$ together
with any two-edge linear forest in $Q[T]$ leaves $|A|+1$ components,
another contradiction. Thus $p=3$, so $|A|=a-1$ and $Q[B]=K_3$.

If $Q[T]$ contained an edge, that edge plus a spanning path of $Q[B]$
would once again leave $|A|+1$ components. Therefore $T$ is independent,
so every vertex of $T$ except possibly one is adjacent to all of $A$,
and the lone exception misses at most one vertex of $A$. Let $b_1b_2b_3$
be a spanning path of $Q[B]$, and write $A=\{x_1,\ldots,x_{a-1}\}$ and
$T=\{t_1,\ldots,t_a\}$. Consider the two vertex-disjoint paths
\[
 b_1b_2b_3x_1t_1x_2t_2\cdots x_{q-1}t_{q-1}
 \quad\text{and}\quad
 t_q x_q t_{q+1} x_{q+1}\cdots x_{a-1}t_a,
\]
which have orders $2q+1$ and $2s+1$. If a missing edge $xt$ exists
between $A$ and $T$, we simply place $t$ in one of the $a\ge3$
$T$-positions not adjacent to $x$ (at most two positions are adjacent to
$x$), and assign the remaining vertices arbitrarily. This gives a
spanning copy of $F$ in $Q$.
\end{proof}

Next, we translate the connected two-odd-path bounds of Yuan and
Zhang~\cite{YZ} and Dong et al.~\cite{DNYZ} into our notation.

\begin{lemma}\label{lem:two-odd}
Let $F=P_{2u+1}\cup P_{2v+1}$ with $u\ge v\ge1$ and $(u,v)\ne(1,1)$, and
put $a=u+v$ and $m=2a+2$. If $G$ is a connected $F$-free graph of order
$n\ge m$, then $e(G)\le\max\{B_m(n),J_a(n)+1\}$.
\end{lemma}
\begin{proof}
Suppose first that $v=1$. When $u=2$, Yuan and
Zhang~\cite[Theorem~1.8]{YZ} showed that
$e(G)\le\max\{\ex(n,P_5),B_8(n),J_3(n)+1\}$, and by the Erd\H{o}s--Gallai
bound we have $\ex(n,P_5)\le3n/2\le2n-2=J_3(n)+1$ for all $n\ge8$.

Now let $u\ge3$, and assume $e(G)\ge B_m(n)$, as otherwise the bound is
immediate. Set $D_u(n)=\binom{2u+1}{2}+n-2u-1$. If $G$ is
$P_{2u+3}$-free, then \eqref{eq:bgls} with $k=2u+2$ yields
$e(G)\le\max\{D_u(n),J_{u+1}(n)+1\}$. If instead $G$ contains
$P_{2u+3}$, Case~2 in the proof of \cite[Theorem~1.8]{YZ} applies
directly (it uses only that $G$ is connected, $(P_{2u+1}\cup P_3)$-free,
contains $P_{2u+3}$, and satisfies $e(G)\ge B_m(n)$): it rules out
$n\le3u+5$ by giving $e(G)<B_m(n)$, and for $n\ge3u+6$ it yields
$e(G)<J_{u+1}(n)+1$. To absorb $D_u(n)$, note that
$D_u(n)\le\binom{2u+3}{2}\le B_m(n)$ for $n\le6u+4$, while at $n=6u+4$
we have $J_{u+1}(n)+1-D_u(n)=(u-1)(7u+4)/2>0$, and this difference grows
with slope $u-1$ for larger $n$.

Suppose now that $v\ge2$. By Dong et al.~\cite[Theorem~4.1]{DNYZ},
\[
 e(G)\le\max\left\{\binom{2a}{2}+n-2a,\,
 \binom{a+1}{2}+(a-1)(n-a-1)\right\}
 =\max\{D_a(n),J_a(n)+1\},
\]
where $D_a(n)=\binom{2a}{2}+n-2a$. For $n\le4a$, we have
$D_a(n)\le\binom{2a+1}{2}\le B_m(n)$; at $n=4a$,
$J_a(n)+1-D_a(n)=(3a^2-9a+2)/2>0$ since $a\ge4$, and the difference
increases with slope $a-2$ thereafter.
\end{proof}

\begin{theorem}\label{thm:at-most-two-odd}
Let $F\notin\{P_3,2P_3\}$ be a nonempty linear forest whose components
all have order at least two, with $O\le2$ odd components. Put
$a=\nu(F)$ and $m=|F|$, and let $\varepsilon=1$ if all components of $F$
are odd and $\varepsilon=0$ otherwise. If $G$ is a connected $F$-free
graph of order $n\ge m$, then
$e(G)\le\max\{B_m(n),J_a(n)+\varepsilon\}$. Moreover, if $F$ contains a
$P_2$ component or if $O\le1$, then $B_m(n)$ may be replaced by
$\binom{m-1}{2}$.
\end{theorem}
\begin{proof}
Assume first that $O\le1$. If $a=1$, then $F=P_2$ and no connected
$F$-free graph of order at least $2$ exists. If $F=P_2\cup P_3$,
applying \eqref{eq:bgls} with $k=4$ shows that $e(G)\le n-1=J_2(n)$,
since the sole equality graph $K_1\vee(K_2\cup I_{n-3})$ contains $F$.

Set $D_m(n)=\binom{m-2}{2}+n-m+2$ and $C_m=\binom{m-1}{2}$. When $F$ has
no $P_2$ component, we have $e(G)\le\max\{D_m(n),J_a(n)+\varepsilon\}$
by \cite[Lemma~2.1]{YZ} if $F$ has at least two components, and by
\eqref{eq:bgls} if $F$ is a single path. When $F$ has a $P_2$ component
(and $F\ne P_2,P_2\cup P_3$), applying \eqref{eq:bgls} with $k=m-1$
yields $e(G)\le\max\{D_m(n),J_a(n)+O\}$. If $O=1$, the equality graph
for the second term contains $F$ (start an even component with the edge
outside the clique and alternate the remaining vertices between the two
parts), so the extra edge cannot occur and again
$e(G)\le\max\{D_m(n),J_a(n)+\varepsilon\}$. At $n_*=2m-4$, we have
$D_m(n_*)=C_m$ and
$2[J_a(n_*)+\varepsilon-C_m]=3(a-1)(a-2)-O(O+1)+2\varepsilon\ge0$; since
the slopes are $1$ and $a-1$, it follows that
$e(G)\le\max\{C_m,J_a(n)+\varepsilon\}$.

Now let $O=2$. If every component is odd, then Lemma~\ref{lem:two-odd}
gives the claim. Otherwise $F$ has an even component, so
$\varepsilon=0$. Put $q=a-1\ge2$, $k=m-1=2q+3$, and $K=\binom{k}{2}$.
Because $P_m$ contains $F$, every longest path in $G$ has order at most
$k$.

\medskip\noindent\emph{Case 1. $G$ is $P_k$-free.}\ 
By \eqref{eq:bgls} with parameter $2q+2$, we have
$e(G)\le\max\{D(n),L_q(n)+1\}$, where $D(n)=\binom{2q+1}{2}+n-2q-1$. At
$n_*=6q+4$, $D(n_*)=K$ and $L_q(n_*)+1-K=(q-1)(7q+4)/2>1$, so comparing
slopes ($1$ versus $q$) gives $D(n)\le\max\{K,L_q(n)\}$ for all $n$. If
$e(G)>\max\{K,L_q(n)\}$, then equality must hold in the second term of
\eqref{eq:bgls}, forcing $G\cong K_q\vee(K_2\cup I_{n-q-2})$. Starting
an even component of $F$ with the edge of $K_2$ and alternating the
remaining vertices between $K_q$ and $I_{n-q-2}$ uses $q=a-1$ vertices
of $K_q$ and $m-q$ vertices outside it, producing a copy of $F$, a
contradiction.

\medskip\noindent\emph{Case 2. $F=H\cup P_2$ and $G$ contains a path of
order $k=|H|+1$.}\ 
By Lemma~\ref{lem:p2-attachment}, the vertices outside this path form an
independent set, and the maximum number $s$ of neighbors of an outside
vertex on the path satisfies $1\le s\le q-2$ (if $q=2$, this case is
vacuous). Applying \eqref{eq:bgls} to the path together with an outside
vertex of degree $s$ gives
\[
 e(G)\le K-\min\{A(s),C(s)\}+s(n-k),\qquad
 A(s)=\frac{s(4q+5-3s)}{2},\quad C(s)=s+\frac{q(q+1)}{2}.
\]
The bounds $K$ and $L_q(n)$ intersect when $n-k=N_0:=(q+2)(q+3)/(2q)$.
Because $s\le q-2$, we have $A(s)/s-N_0\ge3(q-1)/q>0$ and
$C(s)-sN_0\ge6/q>0$ (the latter decreases with $s$). Both linear upper
bounds therefore lie strictly below $\max\{K,L_q(n)\}$ at $n=k+N_0$, and
hence for all $n$.

\medskip\noindent\emph{Case 3. Every component of $F$ has order at
least three, and $G$ contains $P_k$.}\ 
Here $q\ge3$. Choose a longest path $P$ (necessarily of order $k$)
maximizing $d=\max_{y\notin V(P)}d_P(y)$. By
Lemma~\ref{lem:path-completion}, $d\le a-2=q-1$, so
Lemma~\ref{lem:low-attachment} yields
$e(G)\le\max\{K+\lfloor(n-k)/2\rfloor,L_q(n)\}=\max\{B_m(n),J_a(n)\}$.
\end{proof}

\section{Further lemmas for the connected case}\label{sec:coretools}

To handle connected graphs when $H$ has at least three odd components,
we work with the \emph{$j$-core} of a graph ($j\ge0$), defined as its
unique largest induced subgraph of minimum degree at least $j$, which is
obtained by iteratively deleting vertices of current degree less than
$j$. Recall that a graph is \emph{$d$-degenerate} if every nonempty
induced subgraph contains a vertex of degree at most $d$; thus every
subgraph of minimum degree at least $j$ lies inside the $j$-core, and
the $j$-core is empty precisely when the graph is $(j-1)$-degenerate. We
write $\delta(G)$ for the \emph{minimum degree} of $G$, and continue to
use $J_a(x)=(a-1)x-\binom{a}{2}$.

We first record a minimum-degree corollary of a theorem of Egawa and
Ota~\cite{EO}.

\begin{corollary}\label{cor:min-deg-embed}
Let $L$ be a nonempty linear forest whose components have order at least
two, and let $X$ be a connected graph with $|X|\ge|L|$. If
$\delta(X)>\nu(L)$, then $X$ contains $L$. The weaker assumption
$\delta(X)\ge\nu(L)$ also suffices unless $L$ is a union of copies of
$P_3$ or a union of two equal odd paths, and it suffices unconditionally
when $|X|=|L|$.
\end{corollary}
\begin{proof}
Write $L=P_{n_1}\cup\cdots\cup P_{n_k}$. Define
$\sigma_3(X)=\min\{d_X(x)+d_X(y)+d_X(z):\{x,y,z\}\text{ is independent}\}$
if $X$ has an independent set of size $3$, and $\sigma_3(X)=\infty$
otherwise. Because $\sum_i\lfloor n_i/2\rfloor=\nu(L)$, the
minimum-degree hypothesis implies $\sigma_3(X)\ge3\nu(L)$, so Egawa and
Ota~\cite[Theorem~1]{EO} applies.

It remains to check the three exceptional families in \cite{EO}:
\begin{itemize}
\item $L=kP_3$ and $X=K_1\vee(K_{a_1}\cup K_{a_2}\cup K_{a_3})$ with
$a_i\equiv2\pmod{3}$ and $a_1+a_2+a_3=3k$. Here
$\delta(X)=\min\{a_1,a_2,a_3\}\ge\nu(L)=k$, which forces
$a_1=a_2=a_3=k$, so $\delta(X)=\nu(L)$ and $|X|=3k+1>|L|$.
\item $L=2P_{2q+1}$ and $X=K_1\vee mK_{2q}$ with $m\ge3$. Here
$\delta(X)=2q=\nu(L)$ and $|X|=1+2qm>4q+2=|L|$.
\item $L=kP_3$ and $X=mK_2\vee B$ with $|B|=k-1$ and $m\ge k+1$. Every
vertex in $mK_2$ has degree $k$, so $\delta(X)\le k=\nu(L)$ and
$|X|=2m+k-1\ge3k+1>|L|$.
\end{itemize}
In all three families we have $\delta(X)\le\nu(L)$ and $|X|>|L|$, which
yields all assertions of the corollary.
\end{proof}

We also use the following standard consequence of Dirac's
theorem~\cite{Dirac}.

\begin{lemma}\label{lem:dirac-path}
Every connected graph $X$ of order $v$ and minimum degree $d$ contains a
path of order at least $\min\{v,2d+1\}$. In particular, $X$ has a
Hamiltonian path whenever $d\ge(v-1)/2$.
\end{lemma}

\subsection{Deleting an odd path}

Fix $b\ge1$ and set $k=2b+1$. A vertex subset $U\subseteq V(G)$ is
called a \emph{$P_k$-set} if $G[U]$ contains a Hamiltonian path, and we
write $\ell_G(U)=e(G)-e(G-U)$ for the number of edges of $G$ incident
with at least one vertex of $U$. Given a linear order on a vertex set
$X$, let $d_X^+(v)$ be the number of neighbors of $v$ in $X$ appearing
later in the order. For an ordered graph $Q$ and $U\subseteq V(Q)$, let
$e_Q^-(U)$ denote the number of edges directed from $V(Q)\setminus U$
into $U$.

\begin{lemma}\label{lem:ordered-deletion}
Let $k=2b+1$.
\begin{enumerate}
\item[(i)] Let $H$ be a graph of order $k+1$ containing $P_k$. If all
$P_k$-sets of $H$ share a common vertex, then some vertex $v$ with
$P_k\subseteq H-v$ satisfies $d_H(v)\le b$.
\item[(ii)] If $H$ is an ordered graph of order $h>k$ containing $P_k$,
then there exists $v\in V(H)$ such that $H-v$ contains $P_k$ and
$d_H^+(v)\le b$.
\item[(iii)] If $H$ is an ordered graph of order $h\ge k$ containing
$P_k$, then $H$ contains a $P_k$-set $U$ with $e_H^-(U)\le b(h-k)$. The
same bound holds for any acyclic orientation of $H$.
\item[(iv)] In the setting of (iii), if nonnegative vertex weights
satisfy $w(v)+d_H^+(v)\le D$ for all $v\in V(H)$, then some $P_k$-set
$U$ satisfies $\sum_{u\in U}w(u)+\ell_H(U)\le kD+b(h-k)$.
\end{enumerate}
\end{lemma}
\begin{proof}
For (i), assume for a contradiction that every removable vertex (one
whose deletion leaves a $P_k$) has degree at least $b+1$. Pick such a
vertex $x$, and let $q_1\cdots q_k$ be a Hamiltonian path of $H-x$. If
$x$ has two consecutive neighbors on this path, we insert $x$ between
them; otherwise its at least $b+1$ neighbors must be
$q_1,q_3,\ldots,q_{2b+1}$, so $x$ can be appended at either end. In both
cases $H$ has a Hamiltonian path. Its two endpoints are removable and
have degree sum at least $2b+2=|H|$, so closure produces a Hamiltonian
cycle in $H$. Deleting any vertex of this cycle leaves a $P_k$,
contradicting that all $P_k$-sets share a vertex.

We prove (ii) by induction on $h$. For $h=k+1$, part (i) applies if all
$P_k$-sets intersect; otherwise every vertex is removable, and the last
vertex in the order has forward degree $0$. For $h>k+1$, let $x$ be the
first removable vertex in the order. By induction, $H-x$ has a vertex
$u$ with $d_{H-x}^+(u)\le b$ such that $H-x-u$ contains $P_k$. Since $u$
comes after $x$ in the order, restoring $x$ does not change the forward
degree of $u$.

For (iii), iterate (ii) while preserving the induced order until $k$
vertices remain. Each edge entering the resulting $P_k$-set $U$ is
counted precisely once, at the step when its earlier endpoint is
deleted, giving $e_H^-(U)\le b(h-k)$. Any acyclic orientation admits a
topological ordering, yielding the second statement. Finally, for every
$U\subseteq V(H)$,
\begin{equation}\label{eq:weighted-identity}
 \sum_{u\in U}w(u)+\ell_H(U)
 =\sum_{u\in U}\bigl(w(u)+d_H^+(u)\bigr)+e_H^-(U),
\end{equation}
from which (iv) follows immediately via (iii).
\end{proof}

We say that a graph is \emph{hereditarily $b$-sparse} if every connected
subgraph $J$ satisfies $e(J)\le b|J|$.

\begin{theorem}\label{thm:rooted-deletion}
Let $G$ be a connected graph, let $R\ne\varnothing$ induce a connected
subgraph, and set $Y=V(G)\setminus R$ and $N=|Y|$. Suppose that $G[Y]$
contains $P_k$ (where $k=2b+1$) and that $Y$ is linearly ordered so that
$d_R(y)+d_Y^+(y)\le D$ for every $y\in Y$, where $D\ge b$ is an integer.
Set $A_b(D)=\max\{0,k^2-b^2-2b-bD\}=\max\{0,3b^2+2b+1-bD\}$. Then there
exists a $P_k$-set $U\subseteq Y$ such that every component of $G-U$
disjoint from $R$ is hereditarily $b$-sparse and
$\ell_G(U)\le kD+b(N-k)+A_b(D)$. In particular, $A_b(D)=0$ whenever
$b=1$ and $D\ge6$, or $b\ge2$ and $D\ge3b+3$.
\end{theorem}
\begin{proof}
We first establish an unrooted bound for $D$-degenerate graphs: if
$D\ge b$ and a $D$-degenerate graph $H$ of order $h\ge k$ contains
$P_k$, then $H$ has a $P_k$-set $U$ with
\begin{equation}\label{eq:degenerate-bound}
 \ell_H(U)\le T_b(D,h):=bh+(b+1)D-b^2-2b.
\end{equation}
We proceed by induction on $h$. If $h\le D+b$, then any $k$-vertex
subset has at most $\binom{k}{2}+k(h-k)=kh-k(b+1)$ incident edges, and
$T_b(D,h)-(kh-k(b+1))=(b+1)(D-h+2b+1)-b^2-2b\ge1$.

Now let $h\ge D+b+1$, and pick a vertex $v$ of minimum degree $d\le D$.
If $H-v$ is $P_k$-free, the Erd\H{o}s--Gallai bound~\cite{EG} gives
$e(H)\le D+(b-\tfrac{1}{2})(h-1)\le T_b(D,h)$, since the difference is
$b(D-b)+(h-k)/2\ge0$. Suppose instead that $H-v$ contains $P_k$. By
induction, $H-v$ has a $P_k$-set $U$ with
$\ell_{H-v}(U)\le T_b(D,h-1)$. Reintroducing $v$ raises the edge loss by
$d_U(v)$, which is at most $b=T_b(D,h)-T_b(D,h-1)$ unless
$d_U(v)\ge b+1$. Assume $d_U(v)\ge b+1$, and set $W=U\cup\{v\}$ and
$Z=V(H)\setminus W$. If some removable vertex $u\in W\setminus\{v\}$
satisfies $d_Z(u)\ge d-b$, then
$\ell_H(W\setminus\{u\})-\ell_{H-v}(U)=d-d_Z(u)\le b$, proving
\eqref{eq:degenerate-bound}. Otherwise $d_Z(u)\le d-b-1$ for every
removable $u\in W\setminus\{v\}$.

In that case all vertices of $W$ must be removable: if all $P_k$-sets of
$H[W]$ shared a vertex, Lemma~\ref{lem:ordered-deletion}(i) would
produce a removable vertex of degree at most $b$ in $H[W]$, which cannot
be $v$ (as $d_{H[W]}(v)\ge b+1$) and would therefore have degree at most
$b+(d-b-1)=d-1$ in $H$, contradicting the choice of $v$. Setting
$q=d_{H[W]}(v)$ and picking any $u\in W\setminus\{v\}$, we obtain
\begin{align*}
 \ell_H(W\setminus\{u\})
 &\le\left(q+\binom{k}{2}\right)+(d-q)+(k-1)(d-b-1)\\
 &=kd-b\le kD-b\le T_b(D,h),
\end{align*}
where the last step uses $b(h-D-b-1)\ge0$. This proves
\eqref{eq:degenerate-bound}.

Returning to $G$, put $L=kD+b(N-k)$ and $A=A_b(D)$. Applying
Lemma~\ref{lem:ordered-deletion}(iv) to $G[Y]$ with weights $w(y)=d_R(y)$
yields a $P_k$-set in $Y$ of loss at most $L$. Among all $P_k$-sets
$U\subseteq Y$ with $\ell_G(U)\le L+A$, choose one maximizing the order
of the component of $G-U$ containing $R$.

Suppose for a contradiction that some component $C$ of $G-U$ disjoint
from $R$ fails to be hereditarily $b$-sparse. Then $C$ contains a
connected subgraph $J$ with $e(J)>b|J|$, and hence contains $P_k$ by the
Erd\H{o}s--Gallai theorem. Because $G$ is connected, every component of
$G-U$ sends an edge to the connected set $U$, so $G-C$ is connected and
every edge leaving $C$ ends in $U$. Moreover, $G[C]$ is $D$-degenerate
(under the induced order from $Y$) and $|C|\le N-k$. By
\eqref{eq:degenerate-bound}, $C$ contains a $P_k$-set $V$ with
\begin{align*}
 \ell_G(V) &\le b|C|+(b+1)D-b^2-2b+k^2\\
 &\le L+k^2-b^2-2b-bD\le L+A,
\end{align*}
where $k^2$ upper-bounds the edges between $V$ and $U$. In $G-V$, the
connected graph $G-C$ remains intact and contains both $R$ and $U$, so
the component containing $R$ strictly grows, contradicting the
maximality of $U$. The final remark follows directly from the formula
for $A_b(D)$.
\end{proof}

We will also use P\'osa's theorem~\cite{Posa} in its contrapositive
form (see also F\"uredi, Kostochka, and Luo~\cite[Theorem~4]{FKL}): if a
graph on $v\ge3$ vertices has no Hamiltonian cycle, then there exists an
integer $1\le i<v/2$ such that at least $i$ vertices have degree at most
$i$.

The next lemma is a rooted variant of the Erd\H{o}s--Gallai bound.

\begin{lemma}\label{lem:marked-root}
Let $C$ be a connected graph and let $A\subseteq V(C)$ be nonempty. If
no vertex of $A$ is an endpoint of a path of order $h+1$ in $C$ (where
$h\ge1$), then
\[
 e(C)\le\frac{h-1}{2}|C|+\frac{|V(C)\setminus A|}{2}.
\]
\end{lemma}
\begin{proof}
We use induction on $h$. For $h=1$, every vertex of $A$ is isolated, so
connectedness forces $|C|=1$. Let $h\ge2$, and set $q=|C|$ and
$z=|V(C)\setminus A|$. Since $C[A]$ is $P_{h+1}$-free, the
Erd\H{o}s--Gallai theorem~\cite{EG} provides a vertex $x\in A$ with
$d_A(x)\le h-1$, whence $d_C(x)\le h-1+z$.

In every component $D$ of $C-x$, the set $A_D=N_C(x)\cap V(D)$ is
nonempty, and no vertex of $A_D$ can be an endpoint of a $P_h$ in $D$
(prepending $x$ would yield a $P_{h+1}$ starting at $x\in A$). Applying
the induction hypothesis to each component of $C-x$ and summing gives
$e(C-x)\le(h-2)(q-1)/2+(q-1-d_C(x))/2$. Adding $d_C(x)$ and substituting
$d_C(x)\le h-1+z$ completes the proof.
\end{proof}

\subsection{Path factors in a dense core}

For $X\subseteq V(G)$, we abbreviate $e(G[X])$ to $e(X)$, and for
disjoint subsets $A,B\subseteq V(G)$, we write $e(A,B)$ for the number
of edges with one endpoint in $A$ and the other in $B$.

\begin{lemma}\label{lem:rooted-joint}
Let $b\ge1$, let $G$ be a connected graph, and let $V(G)=R\dot\cup Y$
with $R\ne\varnothing$ and $\delta(G[R])\ge b$. If every copy of
$P_{2b+1}$ in $G$ uses at least $b+1$ vertices of $R$, then
$e(R,Y)+e(Y)\le b|R|+(b-1)|Y|$.
\end{lemma}
\begin{proof}
When $b=1$, $G[Y]$ is $P_3$-free, so its components are edges and
isolated vertices. An edge component in $G[Y]$ would have a neighbor in
$R$ by connectedness, forming a $P_3$ with only one vertex in $R$. Thus
$Y$ is independent, and no vertex of $R$ can have two neighbors in $Y$,
so $e(R,Y)\le|R|$.

Now assume $b\ge2$. Delete all edges inside $R$, and then repeatedly
delete any remaining vertex of $R$ of current degree at most $b$, or any
remaining vertex of $Y$ of current degree at most $b-1$, charging the
deleted edges to the removed vertex. Suppose for a contradiction that a
nonempty component $C$ survives. Then $\delta(C)\ge b$, and every vertex
of $C\cap R$ has degree at least $b+1$ in $C$.

Observe that $C$ is $P_{2b+2}$-free: since $R$ is independent in $C$, at
least one of the two consecutive $P_{2b+1}$ subpaths of a $P_{2b+2}$
would contain at most $b$ vertices of $R$. If $|C|\ge2b+1$, then
Lemma~\ref{lem:dirac-path} yields a $P_{2b+1}$ in $C$, which must
alternate between $R$ and $Y$ with both endpoints in $R$. Either
endpoint has at least $b+1$ neighbors in $C$, while at most $b$ of them
can lie on the path, allowing us to extend to a $P_{2b+2}$, a
contradiction. Consequently $q:=|C|\le2b$. Since $\delta(C)\ge b\ge q/2$
and $q\ge b+1\ge3$, Dirac's theorem gives a Hamiltonian cycle in $C$.

If $C$ meets $R$, put $r=|C\cap R|$ and $y=|C\cap Y|$. Because vertices
in $C\cap R$ have degree at least $b+1$ in $C$ and $C\cap R$ is
independent, we have $y\ge b+1$ and hence $r\le b-1$. Open a Hamiltonian
cycle of $C$ at a vertex $x\in C\cap R$, and set $\ell=2b+1-q\ge0$. If
$\ell>0$, choose a path $r_1\cdots r_\ell x$ in $G[R]$ with
$r_1,\ldots,r_\ell\in R\setminus V(C)$ greedily backwards from $x$; at
each step fewer than $r+\ell\le b$ vertices of $R$ are unavailable while
$\delta(G[R])\ge b$. Prepending $r_1\cdots r_\ell$ to the opened cycle
produces a $P_{2b+1}$ with only $r+\ell=2b+1-y\le b$ vertices in $R$, a
contradiction.

If instead $C\subseteq Y$, take a shortest path from $R$ to $C$ (whose
internal vertices lie in $Y\setminus C$) and continue around an opened
Hamiltonian cycle of $C$. Starting at its unique vertex $x\in R$, this
path is followed by at least $q\ge b+1$ consecutive vertices of $Y$.
Attaching a $b$-vertex path in $G[R]$ ending at $x$ and keeping the
first $b+1$ vertices of $Y$ again yields a $P_{2b+1}$ with exactly $b$
vertices in $R$, a contradiction.

Thus the deletion process empties the graph, and summing the charged
degrees gives $e(R,Y)+e(Y)\le b|R|+(b-1)|Y|$.
\end{proof}

Theorem~\ref{thm:ctw} handles the zero-defect case with at least two
components. The next theorem incorporates the ordered one-defect case as
well as the single-path base case needed for induction.

\begin{theorem}\label{thm:dense-core-factor}
Let $H$ be a linear forest whose components have order at least two,
with $a=\nu(H)$ and $O$ odd components. Set $\eta=1$ if all components
of $H$ are odd, and $\eta=0$ otherwise. In the all-odd case, assume
either that $H$ has at least two components with one of order at least
five, or that $H$ is a single path of order at least seven.

Let $G$ be a graph of order $|G|\ge m:=2a+O$ with a partition
$V(G)=R\dot\cup Y$, where $c=|R|\ge a+1+\eta$ and $Y$ is linearly
ordered. Suppose that nonnegative integers $\lambda$ and $\varepsilon_y$
($y\in Y$) satisfy $\binom{c}{2}-e(G[R])\le\lambda$,
$\lambda+\sum_{y\in Y}\varepsilon_y\le\eta$, and
\begin{equation}\label{eq:dense-core-degree}
 d_R(y)+d_Y^+(y)\ge a-1-\varepsilon_y \qquad (y\in Y).
\end{equation}
Then $G$ contains a copy of $H$.
\end{theorem}
\begin{proof}
If $c\ge m$, then $G[R]$ misses at most one edge and hence contains
$P_m$, which contains $H$. Assume $c<m$. Retaining $R$ and the last
$m-c$ vertices of $Y$ preserves all forward degrees of those vertices,
so we may assume $|G|=m$.

Suppose first that $H$ has $t_H\ge2$ components (of which $e_H$ are
even) and that $\sum_{y\in Y}\varepsilon_y=0$. Every vertex of $Y$ has
degree at least $a-1$, while every vertex of $R$ has at least $a$
neighbors in $R$ (since $c\ge a+1$ with $\lambda=0$ when $\eta=0$, and
$c\ge a+2$ with $\lambda\le1$ when $\eta=1$). Hence
$\sigma_2(G)\ge2a-2=m-t_H+e_H-2$, and Theorem~\ref{thm:ctw} gives a
spanning copy of $H$.

We may therefore assume that either $t_H=1$ or
$\sum_{y\in Y}\varepsilon_y=1$. Suppose for a contradiction that $G$ has
no Hamiltonian path, and put $N=|Y|$. We first claim that $N\ge a$.
Assume $N\le a-1$. If an edge is missing inside $R$, then
$\eta=\lambda=1$ and all $\varepsilon_y=0$; the two ends of this
nonedge have degree sum at least $2(c-2)\ge m-1$, so Hamiltonian-path
closure adds the edge. With $R$ now complete, closure adds every missing
edge between $R$ and any $y\in Y$ with $\varepsilon_y=0$ (as their
degree sum is at least $c+a-2\ge m-1$). If a single vertex $y\in Y$ has
$\varepsilon_y=1$, then $\eta=1$ and $a\ge3$, so every vertex of $R$ now
has degree at least $m-2$ while $d(y)\ge a-2\ge1$, and closure joins $y$
to all of $R$ as well. Then $R$ is a universal clique of size $c>N$, so
alternating vertices of $R$ and $Y$ gives a Hamiltonian path, a
contradiction. Thus $N\ge a$.

If $t_H=1$, then $N=m-c\le a-1$ regardless of the parity of $m$, a
contradiction. Hence $t_H\ge2$, which forces $\sum_{y\in Y}\varepsilon_y=1$,
$\eta=1$, and $\lambda=0$, so all components of $H$ are odd and $G[R]$
is complete. We now proceed by induction on $O$.

When $O=2$, we have $a\le N=m-c\le(2a+2)-(a+2)=a$, so $N=a$ and
$c=a+2$. Since $H\ne2P_3$ by hypothesis, Lemma~\ref{lem:two-odd-defect}
(with $S=R$, $T=Y$, and $\delta_y=\varepsilon_y$) yields a spanning copy
of $H$.

Now let $O\ge3$, and let $P_{2b+1}$ be a shortest component of $H$; then
$bO\le a$ and $a\ge2b+2$. Note that $G$ is connected (following forward
neighbors from any vertex of $Y$ eventually reaches $R$, as $a-2\ge1$)
and $\delta(G[R])=c-1\ge a+1\ge b$. If every $P_{2b+1}$ in $G$ used at
least $b+1$ vertices of $R$, then summing \eqref{eq:dense-core-degree}
over $Y$ and applying Lemma~\ref{lem:rooted-joint} would give
\[
 (a-1)N-1\le bc+(b-1)N,\qquad
 a^2\le aN\le b(2a+O)+1\le2ab+a+1,
\]
so $a(a-2b-1)\le1$, contradicting $a\ge2b+2$.

Pick a copy of $P_{2b+1}$ using the minimum number $j\le b$ of vertices
of $R$, and let $S$ be its vertex set in $R$. Place $S$ after $Y$ while
keeping the given order on $Y$. By the minimality of $j$, every
$P_{2b+1}$ in $G[Y\cup S]$ contains $S$. Repeatedly applying
Lemma~\ref{lem:ordered-deletion}(ii) until $2b+1$ vertices remain
deletes only vertices of $Y$; let $P$ be a Hamiltonian path of the
surviving subgraph. When any $y\in Y\setminus V(P)$ was deleted, all
vertices of $S$ and all later vertices of $V(P)\cap Y$ were still
present, so $d_S(y)+d_{V(P)\cap Y}^+(y)\le b$.

Deleting $P$ and keeping the induced order on $Y\setminus V(P)$ leaves a
complete core $R\setminus S$ of size at least $(a-b)+2$, with
\[
 d_{R\setminus S}(y)+d_{Y\setminus V(P)}^+(y)\ge(a-b)-1-\varepsilon_y
 \qquad(y\in Y\setminus V(P))
\]
and total defect at most $1$. The forest $H-P_{2b+1}$ has $O-1$ odd
components, matching number $a-b$, and still contains a component of
order at least $5$ (if $b=1$, a component of order at least $5$ remains;
if $b\ge2$, at least two remain). By induction, $G-P$ contains
$H-P_{2b+1}$, completing the copy of $H$.
\end{proof}

Note that the exclusion of a single $P_5$ in
Theorem~\ref{thm:dense-core-factor} is necessary: $K_4\cup P_1$
satisfies the degree conditions with $a=2$ and one allowed defect, yet
contains no $P_5$.

\subsection{The core obtained by the deletion process}

Throughout this subsection, $H$ is a linear forest with all components
of order at least two that is not a union of copies of $P_3$, having
$O\ge3$ odd components. Put $a=\nu(H)$, $m=2a+O$, $E=a-O\ge1$, and
$K=\binom{m-1}{2}$, and let $\varepsilon=1$ if all components of $H$ are
odd and $\varepsilon=0$ otherwise.

Core and disintegration techniques for linear forests have been used in
stability results by Chen and Zhang~\cite{ChenZhang} and Zhu and
Chen~\cite{ZhuChen}. The proposition below provides an all-order core
structure that preserves both the deletion ordering and the exact edge
surplus over $J_a(n)+\varepsilon$.

\begin{proposition}\label{prop:core}
Let $G$ be a connected $H$-free graph of order $n\ge m$ with
$e(G)>\max\{K,J_a(n)+\varepsilon\}$. Then the $a$-core $R$ of $G$ is
connected. Writing $c=|R|$, we have $c\le m-1$, $\delta(R)\ge a$, and
$c\ge a+2+\varepsilon$. If the vertices of $Y=V(G)\setminus R$ are
listed in their deletion order, then $d_R(y)+d_Y^+(y)\le a-1$ for every
$y\in Y$, and $e(R)-J_a(c)>\varepsilon$. Furthermore, if $P_{2b+1}$ is a
shortest odd component of $H$, then $R-S$ is connected and nonempty for
every $S\subseteq V(R)$ with $|S|\le b$.
\end{proposition}
\begin{proof}
We begin with a simple packing claim: if $p\ge2a-1$ and $q\ge a+1$, then
$P_p\cup P_q$ contains $H$. Indeed, it suffices to select a subcollection
of components of $H$ whose total order lies in $[O+1,a+1]$, placing them
on $P_q$ and the remaining components (of total order at most
$m-(O+1)=2a-1$) on $P_p$.

To find such a subcollection, consider first $E=1$. Then $H$ is either
$P_5\cup(O-1)P_3$ or $P_2\cup OP_3$, whose component-subset sums include
$3j$ and $5+3j$, or $3j$ and $2+3j$; in either case one of the two
consecutive integers $O+1,O+2=a+1$ is attained. Next let $E\ge2$. If
$O=3$ and at least two odd components are $P_3$, their sum is
$6\in[4,a+1]$; otherwise pick a shortest odd component of order at least
$5$, say $3+2e$, so that at least two odd components have excess at
least $e$, giving $E\ge2e$ and $3+2e\le E+3<a+1$. Finally, suppose
$O\ge4$ and assume no subset has total order in $[O+1,a+1]$. Take a
subcollection of maximum total order $L\le O$. Every omitted component
$C$ must then have $|C|\ge a+2-L\ge E+2$, so each omitted odd component
contributes at least $\lceil(E-1)/2\rceil$ to
$E=a-O=\sum_{C\text{ odd}}(|C|-3)/2+\sum_{C\text{ even}}|C|/2$. Because
the chosen set contains at most $\lfloor O/3\rfloor$ odd components, at
least four odd components are omitted when $O\ge5$, yielding
$E\ge4\lceil(E-1)/2\rceil>E$. When $O=4$ and $E\ne3$, at least three odd
components are omitted, giving $E\ge3\lceil(E-1)/2\rceil>E$. In the sole
remaining case $O=4,E=3$, writing the four odd orders as $3+2e_i$ with
$\sum_i e_i\le3$, the two smallest $e_i$ sum to at most $1$, so their
two components have total order in $[6,8]\subseteq[5,8]=[O+1,a+1]$. This
establishes the claim.

Now turn to the $a$-core $R$. If $R=\varnothing$, then $G$ is
$(a-1)$-degenerate and $e(G)\le J_a(n)$, a contradiction. The deletion
order immediately gives $d_R(y)+d_Y^+(y)\le a-1$ for all $y\in Y$, and
summing the deleted degrees yields $e(R)-J_a(c)>\varepsilon$. Every
component of $R$ has order at least $a+1$. By
Corollary~\ref{cor:min-deg-embed}, since $\delta(R)\ge a>\nu(L)$ for
every proper subforest $L$ of $H$, and since $H$ is neither $tP_3$ nor
$2P_{2q+1}$, no component of $R$ can have order at least $m$; hence
every component of $R$ has order at most $m-1$.

Suppose $R$ is disconnected. If some component $C_1$ has order at least
$2a-1$, Lemma~\ref{lem:dirac-path} yields a $P_{2a-1}$ in $C_1$, while a
second component $C_2$ has a $P_{a+1}$, and the claim gives $H$. Thus
every component of $R$ has order at most $2a-2$ and is Hamiltonian by
Dirac's theorem. If $R$ had at least three components, a shortest path
in $G$ between two of them (with internal vertices outside $R$) together
with their Hamiltonian paths would yield a path of order at least
$2a+3$, which with a $P_{a+1}$ from a third component again gives $H$.
Therefore $R$ has exactly two components. Connecting their Hamiltonian
paths via a shortest path in $G$ produces a path of order at least
$c+1$, forcing $c\le m-2$. Setting $s=c-2a$, we have
$2\le s\le O-2\le a-3$ (already impossible for $O=3$). Since both
components have order at least $a+1$, convexity gives
$e(R)-J_a(c)\le\{s(s-1)-a(a-1)+2\}/2<0$, contradicting
$e(R)-J_a(c)>\varepsilon$. Hence $R$ is connected and $c\le m-1$.

Set $t=c-a$, $\lambda=\binom{c}{2}-e(R)$, and
$\Delta=\sum_{y\in Y}(a-1-d_R(y)-d_Y^+(y))$. Counting edges inside $R$
and at each deletion step gives
$e(G)-J_a(n)=\binom{t+1}{2}-\lambda-\Delta$. If $\varepsilon=0$ and
$c=a+1$, the strict excess forces $\lambda=\Delta=0$, and
Theorem~\ref{thm:dense-core-factor} yields $H$. If $\varepsilon=1$, then
$c=a+1$ gives excess at most $1$, while $c=a+2$ forces
$\lambda+\Delta\le1$, so retaining $R$ and the last $m-c$ vertices of
$Y$ again produces $H$ by Theorem~\ref{thm:dense-core-factor}. Thus
$c\ge a+2+\varepsilon$.

It remains to verify the connectivity of $R-S$. Consider first $b=1$,
and suppose $v$ is a cutvertex of $R$. Let the components of $R-v$ have
orders $x_1,\ldots,x_s$. Since $\delta(R)\ge a$, each $x_i\ge a$, so
$s\ge2$ and $t=c-a\ge a+1$. By convexity,
$e(R)\le\sum_{i=1}^s\binom{x_i+1}{2}\le\binom{a+1}{2}+\binom{t}{2}$,
whence
\[
 e(R)-J_a(c)\le a+\frac{t(t-2a+1)}{2}
 \le a(1-E)+\frac{E(E+1)}{2},
\]
using $t\le a+O-1=2a-E-1$. If $E\ge2$, then $a\ge E+2$ and the bound is
at most $(-E^2-E+4)/2<0$, impossible. If $E=1$, the bound equals $1$,
which contradicts $e(R)-J_a(c)>\varepsilon$ when $\varepsilon=1$. The
only remaining possibility is $H=P_2\cup(a-1)P_3$ with $a=O+1\ge3$, in
which equality must hold throughout: since $a+t(t-2a+1)/2$ is strictly
increasing for $t\ge a+1$, we get $t=2a-2$, $c=3a-2$, and $R-v$ has
exactly two components, of orders $a$ and $2a-3$, each forming a clique
together with $v$. Call these two cliques minus $v$ the \emph{private
cliques}.

Because $G$ is connected and $n>c$, some edge $yw$ connects $y\notin R$
to $w\in R$. If $w\ne v$, starting at $y$, traversing the private clique
containing $w$, passing through $v$, and traversing the other private
clique yields a $P_m$ and hence $H$. If $w=v$, note that the private
cliques have orders $a$ and $2a-3$. When $a\equiv0\pmod{3}$, we take $yv$
as $P_2$ and partition both private cliques into copies of $P_3$. When
$a\not\equiv0\pmod{3}$, let $D_1$ be the private clique of order
$1\pmod{3}$ and $D_2$ the other; pick $u\in D_1$, take $yvu$ as a $P_3$
and an edge of $D_2$ as $P_2$, and partition the remaining vertices of
both cliques into copies of $P_3$. This yields $H$ in all cases, so $R$
is $2$-connected.

Finally, suppose $b\ge2$ and some set $S$ of size $s\le b$ separates
$R$. Every component of $R-S$ has at least $a-s+1$ vertices, so
$e(R)-J_a(c)\le\binom{t+1}{2}-(a-s+1)(t-1)$. For $t\ge a-s+2$, this
expression increases with both $t$ and $s$. Substituting $t\le a+O-1$
and $s\le b$, then using $a\ge bO$ (the expression decreases in $a$) and
$O\ge3$ (it decreases in $O$), we obtain
\[
 e(R)-J_a(c)\le\frac{-a^2+(2b+1)a+O^2+(2b-3)O-4b+4}{2}
 \le\frac{-(b^2-1)O^2+(2b^2+3b-3)O-4b+4}{2}.
\]
At $O=3$ this equals $(-3b^2+5b+4)/2$, which is negative for $b\ge3$.
For $b=2$, the only nonnegative case is $O=3,a=6$ (so $H=3P_5$), where
the bound is $1$ and $\varepsilon=1$, again contradicting
$e(R)-J_a(c)>\varepsilon$. Since $|R|\ge a+2>b$, $R-S$ is connected and
nonempty.
\end{proof}

\subsection{Bounds for the core}

Let $r\ge2$, $m=2a+r$, and $k=2b+1$ with $b\ge1$. Let
$\varepsilon\in\{0,1\}$, set $K=\binom{m-1}{2}$, and recall
$B_m(x)=K+\lfloor(x-m+1)/2\rfloor$. In this subsection, $H$ is a linear
forest with matching number $a$, $r$ odd components, shortest odd
component $P_k$, and all even components of order at least $2$, with
$\varepsilon=1$ precisely when all components of $H$ are odd (so
$a\ge rb$, and $a\ge rb+1$ if $\varepsilon=0$).

Let $G$ be a connected $H$-free graph of order $n\ge m$ with connected
$a$-core $R$ of order $c=a+t$, and set $Y=V(G)\setminus R$ and $N=|Y|$.
With $Y$ in its deletion order, put
$\delta_y=a-1-d_R(y)-d_Y^+(y)$, $\Delta=\sum_{y\in Y}\delta_y$, and
$\lambda=\binom{c}{2}-e(R)$; these are nonnegative and satisfy
$e(G)-J_a(n)=\binom{t+1}{2}-\lambda-\Delta$. We examine copies of $P_k$
using at most $b$ vertices of $R$. By Proposition~\ref{prop:core},
deleting any such path leaves a nonempty connected subgraph of $R$, and
we assume this property below.

\begin{lemma}\label{lem:path-loss}
Suppose $G$ contains a copy of $P_k$ using at most $b$ vertices of $R$,
and let $j\le b$ be the minimum number of vertices of $R$ used by such a
path. Then there exists a copy $P$ of $P_k$ using $j$ vertices of $R$
such that
\[
 \ell_G(V(P))\le\mathcal{L}_b(n)-s_j,\qquad
 \mathcal{L}_b(n)=bn+(b+1)a-\frac{3b^2+5b+2}{2},\qquad
 s_j=(b-j)t+\binom{b-j+1}{2}.
\]
Moreover, if $j\ge1$, every component $C$ of $G-P$ disjoint from
$R\setminus V(P)$ satisfies $e(C)\le\eta_j|C|$, where
$\eta_1=b-\tfrac{1}{2}$, $\eta_2=b-1$, and $\eta_j=b-(j-1)/2$ for
$j\ge3$ (in particular, $e(C)\le(b-1)|C|$ whenever $j\ge2$). For
$j\ge3$, the stronger inequality $e(C)\le\eta_j(|C|-1)$ holds.
\end{lemma}
\begin{proof}
Fix a copy of $P_k$ attaining the minimum $j$, and let $S=V(P)\cap R$.
Keep the deletion order on $Y$ and append $S$ at the end in any order,
giving each vertex $v\in Y\cup S$ weight $w(v)=d_{R\setminus S}(v)$. By
the minimality of $j$, every $P_k$ in $G[Y\cup S]$ contains $S$. For
$y\in Y$, we have
$w(y)+d_{Y\cup S}^+(y)=d_R(y)+d_Y^+(y)\le a-1$, while over $S$ the sum
of weights and forward degrees is at most
$e(S,R\setminus S)+e(S)\le j(c-j)+\binom{j}{2}=jc-j(j+1)/2$. By
Lemma~\ref{lem:ordered-deletion}(iii), $G[Y\cup S]$ has a $P_k$-set $U$
with $e_{G[Y\cup S]}^-(U)\le b(N-(k-j))$; let $P$ be a Hamiltonian path
of $G[U]$. Since $S\subseteq U$, identity \eqref{eq:weighted-identity}
counts all edges of $G$ incident with $V(P)$, giving
\[
 \ell_G(V(P))\le(k-j)(a-1)+jc-\frac{j(j+1)}{2}+b(N-k+j)
 =\mathcal{L}_b(n)-s_j.
\]
(For $j=0$, Lemma~\ref{lem:ordered-deletion}(iv) with $w(y)=d_R(y)$
gives the same bound.)

Now let $j\ge1$, and let $C$ be a component of $G-P$ disjoint from
$R\setminus S$. Since $G[Y]$ is $P_k$-free, the Erd\H{o}s--Gallai
theorem gives $e(C)\le(b-\tfrac{1}{2})|C|$, settling $j=1$.

For $j\ge2$, pick a shortest path $r_0z_1\cdots z_\ell x$ from $R$ to
$C$ in $G$, where $r_0\in R$, $x\in C$, and
$z_1,\ldots,z_\ell\notin R\cup C$. Since $\delta(G[R])\ge a$ and
$j-1\le b-1$, we can pick a path of order $j-1$ in $R$ ending at $r_0$
(just $r_0$ when $j=2$). If $C$ had a path of order $k-j+1$ starting at
$x$, concatenating these paths would produce a path of order at least
$k$ using only $j-1$ vertices of $R$, contradicting the minimality of
$j$. Applying Lemma~\ref{lem:marked-root} with $A=\{x\}$ yields
$e(C)\le(k-j)|C|/2-1/2$.

When $j\ge3$, put $h=k-j=2\eta_j$ and $q=|C|$. In every component of
$C-x$, mark the neighbors of $x$; none of them can start a $P_h$ in
$C-x$ (else adding $x$ gives a $P_{h+1}$ starting at $x$). Applying
Lemma~\ref{lem:marked-root} with parameter $h-1$ to each component of
$C-x$ and adding $d_C(x)$ gives
$e(C)\le(h-1)(q-1)/2+d_C(x)/2\le h(q-1)/2=\eta_j(q-1)$.

Finally, let $j=2$. By connectedness of $G$, some vertex $x\in C$ has a
neighbor $u\in V(P)$. If $C$ contained a $P_{2b}$ starting at $x$,
prepending $u$ would yield a $P_{2b+1}$ using at most one vertex of $R$,
contradicting $j=2$. Form $C^+$ by attaching a pendant leaf at $x$; then
$C^+$ is connected and $P_{2b+1}$-free. For $q=|C|\le2b-1$, the trivial
clique bound gives $e(C)\le(b-1)q$. For $q\ge2b$, applying
\eqref{eq:bgls} to $C^+$ and subtracting the pendant edge yields
$e(C)\le\max\{q+2b^2-5b+2,\,(b-1)q-(b-1)(b-2)/2\}\le(b-1)q$, as
$(b-1)q-(q+2b^2-5b+2)=(b-2)(q-2b+1)\ge0$.
\end{proof}

\begin{lemma}\label{lem:residual}
Let $P$ be a copy of $P_k$ using $j$ vertices of $R$ with
$\ell_G(V(P))\le\mathcal{L}_b(n)-s+A$, and let $Q$ be the component of
$G-P$ containing $R\setminus V(P)$. Suppose that every other component
$C$ of $G-P$ satisfies $e(C)\le\eta|C|$, where $|Q|>\eta+1$,
$\tfrac{1}{2}\le\eta<a-b-1$, and $0<b+\eta<a-1$. Put $v=m-k-1$,
$K'=\binom{v}{2}$, $q_0=a+t-j$, and $z=a-2b+r-2$. Set $\xi=1/2$, or
$\xi=0$ if the connected upper bound for $H-P_k$ has the constant term
$K'$ in place of $B_{m-k}$. Using $\binom{x}{2}=x(x-1)/2$ for all
integers $x$, define
\[
 I=(\eta-\xi)
 \begin{cases}
  \displaystyle\frac{\max\bigl\{0,\binom{z+1}{2}-\binom{t-j+1}{2}\bigr\}}{a-1-\xi},
   & q_0\le v,\\[7pt]
  q_0-v, & q_0>v.
 \end{cases}
\]
(Note that the two cases agree at $q_0=v$, where $t-j=z$.) Then
\[
 e(G-P)\le\max\bigl\{K'+\eta(n-m+1)-I,\,J_{a-b}(n-k)+\varepsilon\bigr\}.
\]
If $G-P-Q$ is nonempty of order $w$ and satisfies
$e(G-P-Q)\le\eta w-\theta$, then the first term may be decreased by
$\theta$. Moreover, the second term plus $\ell_G(V(P))$ does not exceed
$J_a(n)+\varepsilon$ whenever $s\ge A$.

Without the reductions $I$ and $\theta$, the minimum gap between
$M_0(x)=\max\{K,J_a(x)+\varepsilon\}$ and the sum of the first term and
$\mathcal{L}_b(x)-s+A$ is $\Gamma_r(b,\eta)+s-A$, attained at
$n_0=(5a+4r-2+f_r)/2$, where $f_r=(r(r+1)-2\varepsilon)/(a-1)$ and
\begin{equation}\label{eq:gamma}
 \Gamma_r(b,\eta)=\frac{b+2-\eta}{2}a+r(1-\eta)
 -\frac{(b+1)(b+2)}{2}-\frac{b+\eta}{2}f_r,
\end{equation}
with slopes $-(b+\eta)$ for $x\le n_0$ and $a-1-b-\eta$ for $x\ge n_0$.
\end{lemma}
\begin{proof}
Write $q=|Q|\ge q_0$. By the core-deletion order,
$e(Q)\le C(q):=(a-1)q-\binom{a}{2}+\binom{t-j+1}{2}$, and for $q\le v$
we also have $e(Q)\le\binom{q}{2}$. Both $\binom{q}{2}-\eta q$ and
$C(q)-\eta q$ are increasing on our range. For $q\ge v+1$, the connected
bound for $H-P_k$ gives either $e(Q)\le J_{a-b}(q)+\varepsilon$ (whose
slope $a-b-1>\eta$ means that adding $e(G-P-Q)\le\eta(n-k-q)$ stays at
most $J_{a-b}(n-k)+\varepsilon$) or
$e(Q)\le\min\{C(q),\,K'+\xi(q-v)\}$.

In the latter case, put $D=K'-C(v)=\binom{z+1}{2}-\binom{t-j+1}{2}$. If
$q_0\le v$ and $D\ge0$, the increasing function $C(q)-\eta q$ meets the
nonincreasing function $K'+\xi(q-v)-\eta q$ at $q=v+D/(a-1-\xi)$, so
their pointwise minimum is maximized there with value
$K'-\eta v-(\eta-\xi)D/(a-1-\xi)$ (and $q<v$ gives no larger value). If
$q_0\le v$ and $D<0$, we simply drop the $C(q)$ cap. If $q_0>v$, the
nonincreasing function is maximized at $q=q_0$, losing
$(\eta-\xi)(q_0-v)$ relative to $q=v$. Adding
$e(G-P-Q)\le\eta(n-k-q)$ (or $\eta(n-k-q)-\theta$ when the sharper
estimate holds) yields the stated bound on $e(G-P)$.

Finally, since $J_a(n)-\mathcal{L}_b(n)=J_{a-b}(n-k)$, the second branch
is covered whenever $s\ge A$. For the first branch, the gap against
$M_0(x)$ has negative slope $-(b+\eta)$ for $x\le n_0$ and positive
slope $a-1-b-\eta$ for $x\ge n_0$, so its minimum occurs at $n_0$, where
direct substitution gives \eqref{eq:gamma}.
\end{proof}

\begin{lemma}\label{lem:positive-core-use}
In the setting of Lemma~\ref{lem:residual}, assume $r\ge3$, $b\ge1$,
$a\ge\max\{rb,r+1\}$ (with $a\ge rb+1$ when $\varepsilon=0$), and
$t\ge2+\varepsilon$, and assume the connected bound for $H-P_k$ holds.
If $G$ contains a copy of $P_k$ using at most $b$ vertices of $R$ and
the minimum number $j$ of vertices of $R$ used by such a path is
positive, then $e(G)\le\max\{K,J_a(n)+\varepsilon\}$.
\end{lemma}
\begin{proof}
Suppose first that $b=1$, so $j=1$ and $a\ge r+1\ge4$. Apply
Lemma~\ref{lem:residual} with $\eta=1/2$ (noting that
$|Q|\ge a+t-1\ge5>\eta+1$). Since $s_1=0=A$, the $J_{a-b}$ term is
covered, while the gap for the first term is
$\Gamma_r(1,1/2)=(5a+2r-12-3f_r)/4$. Its derivative in $a$ is
$5/4+3[r(r+1)-2\varepsilon]/(4(a-1)^2)>0$, and at $a=r+1$ it equals
$r-5/2+3\varepsilon/(2r)>0$. Assume therefore that $b\ge2$.

Pick a copy $P$ of $P_k$ using $j\in[1,b]$ vertices of $R$ as in
Lemma~\ref{lem:path-loss}, with $s_j=(b-j)t+\binom{b-j+1}{2}$ and
$\eta_1=b-\tfrac{1}{2}$, $\eta_2=b-1$, $\eta_j=b-(j-1)/2$ for $j\ge3$;
set $\theta_j=0$ for $j\le2$ and $\theta_j=\eta_j$ for $j\ge3$.

If $G-P$ is connected (no component outside $Q$), we may apply
Lemma~\ref{lem:residual} with $\eta=1/2$ and drop $s_j\ge0$. Since
\[
 \frac{\partial}{\partial a}\Gamma_r(b,1/2)
 =\frac{2b+3}{4}+\frac{(2b+1)[r(r+1)-2\varepsilon]}{4(a-1)^2}>0,
\]
the minimum over $a$ occurs at $a=rb$ (when $\varepsilon=1$) or
$a=rb+1$ (when $\varepsilon=0$). For $\varepsilon=0$ and $a=rb+1$, we
have $d\Gamma_r(b,1/2)/dr=(b+1)^2(2b-1)/(4b)>0$. For $\varepsilon=1$
and $a=rb$, the derivative in $r$ is $N_b(r)/(4(br-1)^2)$, where
$N_b(r)=2b^4r^2+3b^3r^2-4b^3r-6b^2r-2b^2-br^2+3b+2r+3$; here
$N_b'(r)=2(b+1)^2(2b-1)(br-1)>0$ and
$N_b(3)=18b^4+15b^3-20b^2-6b+9>0$ for $b\ge2$ (increasing in $b$ with
$N_2(3)=325$). Thus the minimum occurs at $r=3$, where both choices of
$\varepsilon$ exceed $b^2+3b/4-3/2-1/b>0$ (by $(b-2)(2b+1)/(2b(3b-1))$
and $(2b+3)/4$, respectively).

Now assume $G-P$ has a component outside $Q$. Apply
Lemma~\ref{lem:residual} with $\xi=1/2$, and write
$\Phi_j(t)=\Gamma_r(b,\eta_j)+s_j+\theta_j+I_j(t)$ and $z=a-2b+r-2$.
For $2\le t\le z+j$,
\begin{equation}\label{eq:cap}
 I_j(t)=\frac{\eta_j-1/2}{a-3/2}
 \left\{\binom{z+1}{2}-\binom{t-j+1}{2}\right\}\ge0,
\end{equation}
while $I_j(t)=(\eta_j-1/2)(t-j-z)$ for $t\ge z+j$. Thus $\Phi_j$ is
concave on $[2,z+j]$, and for $t\ge z+j$ its slope is
$(b-j)+(\eta_j-1/2)$, which equals $2b-2$ ($j=1$), $2b-7/2$ ($j=2$), or
$2b-3j/2$ ($j\ge3$), all strictly positive. Hence $\Phi_j(t)$ is
minimized at $t=2$ or $t=z+j$.

Both endpoint values increase with $a$ and $r$. Indeed, at $t=2$, put
$D=a-3/2$ and $h=(2-j)(3-j)\ge0$; up to the positive factor
$\eta_j-1/2$, the term $I_j(2)$ is $\{z(z+1)-h\}/(2D)$, whose
$a$-derivative has numerator $z(2D-z)+D+h>0$ since $2D-z=a+2b-r-1>0$.
At $t=z+j$, $I_j=0$ and $\partial s_j/\partial a=b-j\ge0$. Together with
\[
 \frac{\partial}{\partial a}\Gamma_r(b,\eta_j)
 =\frac{b+2-\eta_j}{2}
 +\frac{(b+\eta_j)[r(r+1)-2\varepsilon]}{2(a-1)^2}>0,
\]
this shows monotonicity in $a$. Along the boundaries $a=rb$ (where
$z=(r-2)(b+1)$ and $df_r/dr\le3/(2b)$) and $a=rb+1,\varepsilon=0$ (where
$z=(r-2)(b+1)+1$ and $df_r/dr=1/b$), $I_j(2)$ increases in $r$, and
substituting into \eqref{eq:gamma} shows that the $r$-derivative of each
endpoint value is at least $b/4+3/(8b)$ for $j=1$, $b/2+1/2+3/(4b)$ for
$j=2$, and $b(j-1)/4+(j-2)/2+3(j-1)/(8b)$ for $j\ge3$. Thus it suffices
to check $r=3$, with $a\ge3b$ and $f_3\le4/b$.

For $j\in\{1,2\}$, evaluation at $t=2$ gives
$\Phi_1(2)\ge(3b^2-10b+4)/(4b)$ and $\Phi_2(2)\ge(b-2)^2/(2b)$. At the
right endpoint, $(r,a,\varepsilon)=(3,3b,1)$ gives
$\Phi_1(z+1)\ge(12b^3-7b^2-21b+4)/(4(3b-1))>0$ and
$\Phi_2(z+2)\ge(b-2)(2b-1)(3b+5)/(2(3b-1))\ge0$, while the mixed case
$(3,3b+1,0)$ is larger by $(12b^3-9b^2+17b-4)/(4b(3b-1))$ and
$(2b-1)(3b^2-3b+4)/(2b(3b-1))$, respectively. Every value is nonnegative
except $\Phi_1(2)$ when $b=2,j=1$: there the value without $I_1(2)$ is
$-1/2$, while $I_1(2)=10/9$, giving $\Phi_1(2)=11/18>0$ (and
$\Phi_1(z+1)=3/2$).

Finally, let $j\ge3$ and set $q=b-j\ge0$. At $t=2$,
$I_j(2)\ge(b+q)(q+4)(2b-q-1)/(6(2b-1))\ge b(q+4)/6$, so completing the
square gives
\[
 \Phi_j(2)\ge\frac{1}{2}\left(q-\frac{7b}{12}\right)^2
 +\frac{23b^2}{288}+\frac{5b}{12}-2-\frac{1}{b}
 +\left(\frac{3}{2}-\frac{1}{b}\right)q,
\]
which is positive for $b\ge4$ (at $b=4,q=0$, the non-square part is
$25/36$); for $b=3$ we have $q=0$ and $\Phi_3(2)=7/6$. At $t=z+j$,
\[
 \Phi_j(z+j)\ge\frac{b(b-1)}{4}-2-\frac{1}{b}
 +q\left(\frac{5b}{4}+\frac{1}{2}-\frac{1}{b}\right)-\frac{q^2}{2},
\]
which is positive for $b\ge4$ since $q^2\le bq$ (at $b=4,q=0$ it equals
$3/4$). For $b=j=3$, the value is nonnegative at every allowed point
except the real intersection point for $(r,a,\varepsilon)=(3,9,1)$:
indeed, $(3,10,0)$ gives $2/3$ and $(4,12,0)$ gives $16/11$ (and valid
mixed cases with $r=4$ have $a\ge13$). At $(r,a,\varepsilon)=(3,9,1)$,
the real minimum is $-5/8$ at $n_0=225/8$, with slopes $-5$ and $3$, so
the two nearest integers $n=28$ and $n=29$ give gaps $0$ and $2$,
respectively. Thus no integer $n$ produces a strict excess.
\end{proof}

\section{Proof of the connected theorem}\label{sec:connected}

P\'osa-type degree conditions for embedding forests were investigated by
Babu and Diwan~\cite{BabuDiwan}. To finish the connected case, we need a
sharper structural description of a non-traceable graph containing a
dense core.

\begin{proposition}\label{prop:posa-structure}
Let $a\ge1$, and let $Q$ be a graph of order $m=2a+r$ with
$V(Q)=R\dot\cup U$, where $|R|=a+t$ and $1\le|U|=d=a+r-t\le a$. Assume
that $\delta(Q[R])\ge a$ and $d_Q(u)\ge d$ for every $u\in U$. Set
$\lambda=\binom{|R|}{2}-e(R)$, and suppose
$\lambda+\Delta\le\binom{t+1}{2}-\varepsilon-1$ for some
$\Delta,\varepsilon\ge0$. If $Q$ has no Hamiltonian path, then there
exist an integer $s$ with $1\le s\le\lfloor r/2\rfloor$ and a partition
$R=L\dot\cup S$ such that
\[
 |L|=t-r+s,\qquad |S|=a+r-s,\qquad a\le d_R(x)\le a+s-1\quad(x\in L).
\]
Putting $\ell=|L|$, $H_0=\sum_{x\in L}(a+s-1-d_R(x))$, and
$\lambda_S=\binom{|S|}{2}-e(S)$, we have
\begin{equation}\label{eq:posa-count}
 \lambda=\binom{\ell}{2}+\ell(r-2s+1)+H_0+e(L)+\lambda_S.
\end{equation}
In particular, writing $c=|R|$,
\begin{equation}\label{eq:posa-endpoints}
 \lambda\ge\min\left\{
 \frac{(a-d)(2c+d-3a+1)}{2},\,
 \frac{(c-d)(c-d+2)}{8}\right\}.
\end{equation}
Moreover:
\begin{enumerate}
\item[(i)] If $r=2$ and $t\ge7$, then the Hamiltonian-path closure of
$Q$ contains a set $A\subseteq R$ of $a$ universal vertices, and
$V(Q)\setminus A$ is independent in $Q$.
\item[(ii)] If $r=3$ and $t\ge8$, then the Hamiltonian-path closure of
$Q$ contains a set $A\subseteq R$ of $a$ universal vertices, and
$e(Q[V(Q)\setminus A])\le1$.
\item[(iii)] If $r=4$ and $t\ge14$, then $Q$ contains a set $W$ of
order $a+3$ with $e(Q[W])\le1$.
\end{enumerate}
\end{proposition}
\begin{proof}
We first note a quick linear-forest bound. Let $J$ be a graph of maximum
degree $\Delta(J)\le h+2$ ($h\ge1$), and let $C\subseteq V(J)$ be a set
of vertices of degree at least $h$ in $J$. Then a maximum-edge spanning
linear forest of $J$ has at least $|C|/2$ edges, and strictly more than
$|C|/2$ edges when $C\ne\varnothing$, $h=2$, and $\Delta(J)\le3$.
Indeed, let such a forest have $E$ edges, $p$ nontrivial path
components, and $z$ isolated vertices lying in $C$. By maximality, every
neighbor of an isolated vertex is an internal vertex of the forest.
Counting unused incidences at internal vertices gives $hz\le h(E-p)$, so
$|C|\le z+E+p\le2E$. When $h=2$ and $\Delta(J)\le3$, each internal
vertex has only one unused incidence, giving $2z\le E-p$ and
$|C|\le(3E+p)/2\le2E$; equality would force $p=E$, $z=0$, and all
nonisolated vertices to form a matching inside $C$, which is impossible
because every vertex of $C$ has degree at least $2$ in $J$, allowing the
forest to be extended.

Now adjoin a new universal vertex to $Q$. Since $Q$ has no Hamiltonian
path, the enlarged graph is not Hamiltonian, so P\'osa's theorem yields
an integer $i<(m+1)/2$ such that at least $i$ vertices have degree at
most $i$ in the enlarged graph. In this graph, vertices of $R$ have
degree at least $a+1$, vertices of $U$ have degree at least $d+1$, and
the new vertex has degree $m$. Hence $i>d$, and $i\le a$ is impossible
because only the $d<i$ vertices of $U$ could qualify. Thus $i=a+s$ for
some $1\le s\le\lfloor r/2\rfloor$.

At least $i-d=t-r+s$ of these low-degree vertices belong to $R$. Pick
any $\ell=t-r+s$ of them to form $L$, and set $S=R\setminus L$. Since
$a\le d_R(x)\le a+s-1$ for all $x\in L$, we have
$2e(L)+e(L,S)=\ell(a+s-1)-H_0$, and subtracting $e(S)+e(L,S)+e(L)$ from
$\binom{|R|}{2}$ gives \eqref{eq:posa-count}. Dropping $H_0$, $e(L)$,
and $\lambda_S$ yields $\lambda\ge(i-d)(2c+d-3i+1)/2$, a concave
quadratic in $i$ whose values at the endpoints $i=a$ and $i=m/2$ are the
two terms in \eqref{eq:posa-endpoints}.

For $r\in\{2,3,4\}$, substituting \eqref{eq:posa-count} into the upper
bound on $\lambda+\Delta$ gives:
\begin{align}
 r\in\{2,3\},\ s=1:\quad
 &\lambda=\binom{t}{2}-r+2+e(L)+\lambda_S,
 && \Delta+e(L)+\lambda_S\le t+r-3-\varepsilon, \label{eq:budget23}\\
 r=4,\ s=2:\quad
 &\lambda=\binom{t-1}{2}+H_0+e(L)+\lambda_S,
 && \Delta+H_0+e(L)+\lambda_S\le2t-\varepsilon-2, \label{eq:budget42}\\
 r=4,\ s=1:\quad
 &\lambda=\binom{t}{2}-3+e(L)+\lambda_S,
 && \Delta+e(L)+\lambda_S\le t+2-\varepsilon. \label{eq:budget41}
\end{align}

Next, perform the following closure operation starting from $Q$:
whenever a nonuniversal vertex $x$ reaches current degree at least
$a+r-1$, first add all missing edges from $x$ to $R$ (valid because
$d(x)+d(y)\ge(a+r-1)+a=m-1$ for every $y\in R$), and then add all
missing edges from $x$ to $U$ (valid because after the first step
$d(x)\ge|R|-1$ if $x\in R$ and $d(x)\ge|R|$ if $x\in U$, so
$d(x)+d(u)\ge m-1$ for every $u\in U$). Each such vertex $x$ thus
becomes universal. Stop when no such vertex remains, let $A$ be the set
of all universal vertices in the resulting graph (including those
already universal in $Q$), and set $W=V(Q)\setminus A$. Every added edge
is incident with $A$, so $Q[W]$ is unchanged and every vertex of $W$ has
final degree at most $a+r-2$.

Put $p=|S\setminus A|$, and for $x\in S$ let
$\mu_S(x)=|S|-1-d_S(x)$. Every $x\in S\setminus A$ satisfies
$d_L(x)\le s-1+\mu_S(x)$ (otherwise $d_R(x)>a+r-2$). In $Q[R]$, the
number of nonedges between $S$ and $L$ is $\ell(r-2s+1)+H_0+2e(L)$,
while $A\cap S$ contributes at most $(|S|-p)\ell$ edges to $L$. Summing
the bound on $d_L(x)$ over $S\setminus A$ and using
$\sum_{x\in S}\mu_S(x)=2\lambda_S$ yields
\begin{equation}\label{eq:universal-count}
 (p-r+2s-1)\ell\le p(s-1)+H_0+2e(L)+2\lambda_S.
\end{equation}
When $r\in\{2,3\}$, subtracting $2\ell$ from \eqref{eq:universal-count}
and applying \eqref{eq:budget23} gives
$(p-r-1)\ell\le4r-8-2\varepsilon-2\Delta$; since $\ell=t-r+1\ge6$, this
forces $p\le r+1$ and $|A|\ge a-2$. When $r=4$ and $s=2$, subtracting
$4\ell$ and using \eqref{eq:budget42} gives
$(p-5)\ell\le p+4-2\varepsilon-2\Delta-H_0$, which with $\ell=t-2\ge12$
forces $p\le5$ and $|A|\ge a-3$. When $r=4$ and $s=1$,
\eqref{eq:budget41} gives $(p-5)\ell\le10-2\varepsilon-2\Delta$, and
$\ell=t-3\ge11$ again forces $p\le5$ and $|A|\ge a-2$.

We claim that $|A|\ge a$. Suppose $|A|=a-h$ with $h\ge1$. Then
$\Delta(Q[W])\le h+r-2$, and all $|R|-|A|=t+h$ vertices of $R\cap W$
have degree at least $h$ in $Q[W]$. Let $E$ be the maximum number of
edges in a spanning linear forest of $Q[W]$. Since $Q$ has no
Hamiltonian path, the closure cannot link the components of this forest
through $A$, so the forest must have more than $|A|+1$ components,
giving $E\le|W|-|A|-2=2h+r-2$. Applying our initial linear-forest bound
to $C=R\cap W$ in $Q[W]$ gives a contradiction in every case:
$t+h\le2E\le4h$ ($h\le2$) contradicts $t\ge7$ for $r=2$;
$t+h\le2E\le4h+4$ ($h\le3$) contradicts $t\ge14$ for $r=4$; and for
$r=3$, $h=1$ gives $t+1\le6$ while $h=2$ (using the strict form) gives
$t+2<10$, both contradicting $t\ge8$. Hence $|A|\ge a$.

If $r\in\{2,3\}$, then $|A|\ge a+1$ would leave $|W|\le|A|+1$ singletons
and yield a Hamiltonian path, so $|A|=a$. Likewise, $Q[W]$ can have no
edge when $r=2$ and at most one edge when $r=3$ (two edges would form a
two-edge linear forest with at most $|A|+1$ components). Furthermore,
$A\subseteq R$: if $A\cap U\ne\varnothing$, then $|A\cap R|\le a-1$, and
picking a vertex $x\in R\cap W$ isolated in $Q[W]$ (which exists since
$|R\cap W|\ge t+1\ge9$ when $r=3$, and $Q[W]$ is independent when $r=2$)
gives $d_R(x)\le|A\cap R|\le a-1$, contradicting $\delta(Q[R])\ge a$.
This proves (i) and (ii).

Finally, let $r=4$. We cannot have $|A|\ge a+2$. If $|A|=a+1$, then
$|W|=a+3$ and $W$ is independent (any edge in $Q[W]$ would leave at most
$|A|+1$ path components). If $|A|=a$, then $|W|=a+4$,
$\Delta(Q[W])\le2$, and $Q[W]$ has no three-edge linear forest, so its
nonisolated part is one of $P_2$, $2P_2$, $P_3$, or $K_3$; deleting one
suitable vertex leaves an $(a+3)$-subset spanning at most one edge,
proving (iii).
\end{proof}

\subsection{Path factors in the original graph}

We will make repeated use of two routine consequences of Hall's theorem.
\begin{enumerate}
\item[(H1)] Let $M\ge6$. Suppose $M$ positions are to be filled by $M$
distinct vertices such that each vertex is forbidden from at most $2$
positions, except that at most two vertices may be forbidden from at
most $4$ positions, and every position accepts at least $2$ vertices.
Then a valid assignment exists.
\item[(H2)] Suppose $M$ positions are to be filled by $M$ distinct
vertices such that each vertex is forbidden from at most $q$ positions.
If either every position accepts at least $q$ vertices, or there is a
fixed set of at least $q$ vertices acceptable at every position, then a
valid assignment exists.
\end{enumerate}
Both follow immediately from Hall's condition for any nonempty set
$\mathcal{S}$ of positions: in (H1), the union of available vertices has
size at least $2$ for $|\mathcal{S}|\le2$, at least $M-2\ge|\mathcal{S}|$
for $3\le|\mathcal{S}|\le M-2$, and $M$ for $|\mathcal{S}|\ge M-1>4$; in
(H2), the union has size at least $q$ for $|\mathcal{S}|\le q$ and $M$
for $|\mathcal{S}|>q$.

\begin{theorem}\label{thm:factor-from-W}
Let $F$ be a linear forest with $r\in\{3,4\}$ odd components of orders
at least $2b+1\ge5$, even components of orders at least two, and
matching number $a$. Assume $1\le d\le b$ if $r=3$, and $b\ge3$ with
$1\le d\le\lfloor b/2\rfloor$ if $r=4$. Let $Q$ be a graph with
$V(Q)=R\dot\cup U$, where $|R|=2a+r-d$, $|U|=d$, $\delta(Q[R])\ge a$,
and $d_Q(u)\ge\kappa:=\lceil(a+d+2r-5)/2\rceil$ for all $u\in U$. If $Q$
contains a set $W$ of order $a+3$ with $e(Q[W])\le1$, then $Q$ has a
spanning copy of $F$.
\end{theorem}
\begin{proof}
Set $\tau=r-3$ and $A=V(Q)\setminus W$, so $|A|=a+\tau$. Since
$|W\cap R|\ge a+3-d$, we have $|W\cap R|\ge7$ when $r=3$ and
$|W\cap R|\ge14$ when $r=4$. Because $e(Q[W])\le1$, some vertex
$v\in W\cap R$ is isolated in $Q[W]$; its at least $a$ neighbors in $R$
all lie in $A\cap R$, so $|A\cap R|\ge a$ and $|A\cap U|\le\tau$.

We construct the components of $F$ using alternating templates whose
\emph{$X$-positions} are filled by vertices of $A$ and whose
\emph{$Z$-positions} are filled by vertices of $W$ (with some
$Z$-positions pre-filled by $U$): an odd path $P_{2h+1}$ takes either
$Z_0X_1Z_1\cdots X_hZ_h$ ($h$ $X$-positions, $h+1$ $Z$-positions) or
$X_0Z_1X_1\cdots Z_hX_h$ ($h+1$ $X$-positions, $h$ $Z$-positions), while
an even path $P_{2h}$ takes $Z_0X_1\cdots Z_{h-1}X_h$. Replacing an
adjacent pair $XZ$ by $XZ'X'Z$ extends any initial template using at
most $b$ vertices of $A$ to any prescribed odd order at least $2b+1$,
keeping every unfilled $Z$-position adjacent only to $X$-positions.

\medskip\noindent\emph{Case 1. $A\subseteq R$.}\ 
Set $B=R\setminus A$, $M=|B|=a+3-d$, and $\mu(x)=|B\setminus N_Q(x)|$
for $x\in A$. For every $v\in B$, $\delta(Q[R])\ge a$ gives
$|A\setminus N_Q(v)|\le\tau+d_B(v)$; summing over $B$ and using
$e(Q[W])\le1$ gives
\begin{equation}\label{eq:missing-sum}
 \sum_{x\in A}\mu(x)\le\tau M+2,
\end{equation}
and for every $x\in A$,
\begin{equation}\label{eq:internal-degree}
 d_{Q[A]}(x)\ge a-(M-\mu(x))=\mu(x)+d-3.
\end{equation}
If some pair $x,y\in A$ satisfies $\mu(x)+\mu(y)\ge M-1$, then
\eqref{eq:missing-sum} forces $\tau=1$; pick $z\in A$ maximizing
$\mu(z)$. Since $M\ge14$, we have $\mu(z)\ge\lceil(M-1)/2\rceil\ge7$,
$\sum_{x\in A\setminus\{z\}}\mu(x)\le M-2$, and
$d_{Q[A]}(z)\ge d+2$ by \eqref{eq:internal-degree}. If no such pair
exists, we leave $z$ unselected, and $\mu(x)+\mu(y)\le M-2$ for all
$x,y\in A$.

List $U=\{u_1,\ldots,u_d\}$. Each $u_i$ has at most one neighbor in $W$,
so $|N_A(u_i)|\ge\kappa-1$, and for $1\le i<d$,
\[
 |N_A(u_i)\cap N_A(u_{i+1})|\ge2(\kappa-1)-|A|
 \ge\begin{cases}d-1, & r=3,\\ d, & r=4.\end{cases}
\]
Setting $\mathcal{C}_i=(N_A(u_i)\cap N_A(u_{i+1}))\setminus\{z\}$
(deleting $z$ only if it was selected), each of the $d-1$ sets
$\mathcal{C}_i$ has size at least $d-1$, so Hall's theorem gives
distinct $c_i\in\mathcal{C}_i$. Since $\kappa-d\ge1$ for $r=3$ and
$\kappa-d-1\ge1$ for $r=4$, we can also choose
$c_d\in N_A(u_d)\setminus\{c_1,\ldots,c_{d-1},z\}$. Appending an
unfilled $Z$-position $Z_\ast$ yields the initial template
\begin{equation}\label{eq:outside-path}
 u_1c_1u_2c_2\cdots u_dc_dZ_\ast,
\end{equation}
which absorbs all of $U$ using $d\le b$ vertices of $A$ and therefore
extends to one prescribed odd component.

If $z$ was selected, pick two unused neighbors $v_1,v_2\in N_A(z)$
(possible since $d_{Q[A]}(z)\ge d+2$) and start a second odd component
with $Z_1v_1zv_2Z_2$, so that $z$ is not adjacent to any unfilled
$Z$-position. If $\tau=1$ and no $z$ was selected, give a second odd
component one more $X$-position than $Z$-positions. Give every remaining
odd component one more $Z$-position than $X$-positions, and every even
component equal numbers of each. The total counts of $X$- and
$Z$-positions are then $a+\tau=|A|$ and $a+3=|W|$.

Assign the unused vertices of $A$ to the remaining $X$-positions. There
remain $M=|B|\ge7$ unfilled $Z$-positions, none adjacent to $z$. Any
unfilled $Z$-position has at most two adjacent $A$-vertices $x,y$, which
forbid at most $\mu(x)+\mu(y)\le M-2$ vertices of $B$, so every position
accepts at least $2$ vertices of $B$. Conversely, each vertex of $B$
misses at most $\tau$ vertices of $A$ (at most $\tau+1$ for the at most
two endpoints of the possible edge in $Q[B]$), and each $X$-position
touches at most two unfilled $Z$-positions; thus every ordinary vertex
of $B$ is forbidden from at most $2$ positions and the at most two
exceptions from at most $4$. Applying (H1) completes the spanning copy
of $F$.

\medskip\noindent\emph{Case 2. $A\cap U\ne\varnothing$.}\ 
Here $r=4$, so we write $A=C\dot\cup\{y\}$ with $y\in U$ and $|C|=a$,
and put $X=U\setminus\{y\}$ and $B=R\setminus C$. Every $v\in B$
satisfies $|C\setminus N_Q(v)|\le d_B(v)\le1$, so all vertices of $B$
except at most two are adjacent to all of $C$, and each exception misses
just one vertex of $C$. For every $x\in X$, at most one neighbor lies in
$W$ and at most one is $y$, giving
\begin{equation}\label{eq:X-degree}
 |N_C(x)|\ge\kappa-2
\end{equation}
and, for distinct $x,x'\in X$,
\begin{equation}\label{eq:X-common}
 |N_C(x)\cap N_C(x')|\ge2(\kappa-2)-a\ge d-1.
\end{equation}

If $d=1$, pick any neighbor $w$ of $y$. If $w\in C$, then
$d_R(w)\ge a$ implies $w$ has a neighbor $z\in B$, and we start with
$ywz$; if $w\in B$, then $|C\setminus N_Q(w)|\le1$ gives $c\in N_C(w)$,
and we start with $ywc$.

Suppose $d\ge2$. Pick a neighbor $w$ of $y$, and list
$X=\{x_1,\ldots,x_{d-1}\}$, setting $x_1=w$ if $w\in X$. By
\eqref{eq:X-common} and Hall's theorem, we can pick distinct common
neighbors $c_i\in N_C(x_i)\cap N_C(x_{i+1})$ to form
$P=x_1c_1x_2c_2\cdots c_{d-2}x_{d-1}$, avoiding any one prescribed vertex
of $C$ if desired (since $d-2$ choices remain in each set). Using
\eqref{eq:X-degree}, $a\ge4b$, and $d\le b/2$, we start one odd
component as follows:
\begin{itemize}
\item If $w\in C$, build $P$ avoiding $w$, choose distinct unused
vertices $v\in N_C(x_1)$ and $u\in N_C(x_{d-1})$ (using
$\kappa-d-2\ge1$), and start with $ywZ_\ast vPuZ'_\ast$ (using
$d+1\le b$ vertices of $C$).
\item If $w\in B$, pick $v\in N_C(w)\cap N_C(x_1)$ (size at least
$\kappa-3$), build $P$ avoiding $v$, pick an unused $u\in N_C(x_{d-1})$,
and start with $ywvPu$ (using $d$ vertices of $C$ and $w\in B$).
\item If $w\in X$, take $x_1=w$, build $P$, pick an unused
$u\in N_C(x_{d-1})$, and start with $yPu$ (using $d-1$ vertices of
$C$).
\end{itemize}
Extend this initial template to one odd component, give the other three
odd components one more $Z$-position than $C$-positions, and give every
even component equal numbers of each; the totals are $a=|C|$ and $a+4$.
After filling the remaining $C$-positions arbitrarily from $C$, the
remaining $M\ge14$ vertices of $B$ match the $M$ unfilled $Z$-positions:
all but at most two vertices of $B$ can fill any position, and each
exception is forbidden from at most $2$ positions. Applying (H2) with
$q=2$ completes the spanning copy of $F$.
\end{proof}

Path partitions with prescribed endpoints were studied by Enomoto and
Ota~\cite{EnomotoOta} and Kawarabayashi~\cite{Kawarabayashi}. We need
the following near-complete-core variant.

\begin{lemma}\label{lem:rooted-completion}
Let $a\ge6$ and $s\in\{1,2\}$, and let $R$ be a graph of order $2a+s$
with $\delta(R)\ge a$ and
$\binom{2a+s}{2}-e(R)\le\binom{a+s+1}{2}-1$.
\begin{enumerate}
\item[(i)] For any $x\in V(R)$, either $R$ has a Hamiltonian path
starting at $x$, or $V(R)=A\dot\cup B$ with $|A|=a$, $|B|=a+s$,
$e(R[B])\le s-1$, every vertex of $B$ missing at most $s-1$ vertices of
$A$, and all but at most two vertices of $B$ adjacent to all of $A$.
Moreover, $x\in A$ when $s=1$.
\item[(ii)] Let $F$ be a linear forest with $r\in\{3,4\}$ odd components
of orders at least five, even components of orders at least two, and
matching number $a$, where $(r,s)\in\{(3,1),(4,1),(4,2)\}$. If a new
path $z_1\cdots z_{r-s}$ is attached to $R$ by the single edge
$z_{r-s}x$ for a prescribed $x\in V(R)$, then the resulting graph has a
spanning copy of $F$.
\end{enumerate}
\end{lemma}
\begin{proof}
For (i), attach a pendant leaf $z$ to $x$. If the new graph has a
Hamiltonian path, deleting $z$ leaves a Hamiltonian path of $R$ starting
at $x$. Otherwise apply Proposition~\ref{prop:posa-structure} with
$(r,t,d,\Delta,\varepsilon)=(s+1,a+s,1,0,0)$ to obtain $A\subseteq V(R)$
of size $a$ such that the complement (including $z$) spans at most $s-1$
edges. Setting $B=V(R)\setminus A$, each $v\in B$ has
$d_A(v)\ge a-d_B(v)\ge a-(s-1)$, and since $2e(R[B])\le2$, at most two
vertices of $B$ have neighbors in $B$. When $s=1$, $B\cup\{z\}$ is
independent, so $xz\in E$ forces $x\in A$.

For (ii), if $R$ has a Hamiltonian path starting at $x$, prepending
$z_1\cdots z_{r-s}$ yields a spanning path of order $2a+r$, which splits
into $F$. Otherwise take the partition $V(R)=A\dot\cup B$ from (i), and
let $Z$ denote an unfilled $B$-position. Start one odd component with
$z_1z_2x$ when $r-s=2$ and $x\in A$ (covering $(r,s)=(3,1)$ and
$(4,2)$ with $x\in A$), with $z_1z_2z_3xZ$ when $r-s=3$ (so $s=1$ and
$x\in A$), and with $z_1z_2xyZ$ for some $y\in N_R(x)\cap A$ when
$(r,s)=(4,2)$ and $x\in B$ (possible since $x$ misses at most one vertex
of $A$).

Extend this initial path by alternating unused vertices of $A$ and
unfilled $B$-positions. If the initial path has three more vertices
outside $A$ than inside $A$, give one other odd component one more
$A$-position than $B$-positions, every remaining odd component one more
$B$-position, and every even component equal numbers of each; if it has
only one more vertex outside $A$, give every remaining odd component one
more $B$-position. In all cases $a$ positions are filled from $A$, and
$M=a+s-\theta\ge a\ge6$ unfilled $B$-positions remain (where $\theta=1$
if $x\in B$ was already used, and $\theta=0$ otherwise). All but at most
two remaining vertices of $B$ can fill every position, and each
exception is forbidden from at most $2$ positions, so (H2) with $q=2$
completes the spanning copy of $F$.
\end{proof}

\subsection{Paths using few vertices of the core}

\begin{proposition}\label{prop:no-low-core-path}
Let $H$ be a linear forest with $r\ge3$ odd components, no isolated
vertices, and shortest odd component $P_{2b+1}$, and suppose $H$ is not
a union of copies of $P_3$. Put $a=\nu(H)$, $m=2a+r$, and
$K=\binom{m-1}{2}$, and let $\varepsilon=1$ if all components of $H$ are
odd and $\varepsilon=0$ otherwise. Let $G$ be a connected $H$-free graph
of order $n\ge m$ with $e(G)>\max\{K,J_a(n)+\varepsilon\}$. Assume that
the connected bound for $H-P_{2b+1}$ holds with terms $B_{m-2b-1}$ and
$J_{a-b}+\varepsilon$ (and with $\binom{m-2b-2}{2}$ in place of
$B_{m-2b-1}$ if $H$ has a $P_2$ component). Then no copy of $P_{2b+1}$
in $G$ uses at most $b$ vertices of the $a$-core $R$ from
Proposition~\ref{prop:core}.
\end{proposition}
\begin{proof}
Write $|R|=a+t$. By Proposition~\ref{prop:core}, $t\ge2+\varepsilon$
and deleting at most $b$ vertices from $R$ leaves a connected subgraph
of order at least $a+t-b\ge2b+2$. Also $a>2b+1$ (as $a\ge rb$ and
$H\ne rP_3$). Every value of $\eta$ used below therefore satisfies
$|Q|>\eta+1$, $1/2\le\eta<a-b-1$, and $0<b+\eta<a-1$ in
Lemma~\ref{lem:residual}.

Assume for a contradiction that $G$ has a $P_{2b+1}$ using at most $b$
vertices of $R$, and let $j$ be the minimum number of vertices of $R$
used by such a path. Lemma~\ref{lem:positive-core-use} rules out $j>0$,
so $j=0$, meaning that $G-R$ contains $P_{2b+1}$.

Apply Theorem~\ref{thm:rooted-deletion} with parameter $b$ and $D=a-1$,
which gives error term $A=\max\{0,3b^2+3b+1-ba\}$. Put
$f=(r(r+1)-2\varepsilon)/(a-1)$ and $s_0=bt+\binom{b+1}{2}$. By
Lemma~\ref{lem:residual}, the minimum gap against $K$ (before $I$) is
\[
 \Psi:=\Gamma_r(b,b)+s_0-A=a+bt+r-(r+1)b-1-bf-A,
\]
while the gap against the join term is $s_0-A$. For fixed $r$ and $b$,
$\Psi(a+1)-\Psi(a)=1+b[r(r+1)-2\varepsilon]/(a(a-1))+A(a)-A(a+1)>0$;
increasing $t$ by $1$ adds $b$ to $\Psi$, and changing $\varepsilon$
from $0$ to $1$ adds $2b/(a-1)$, without decreasing $s_0-A$. We may
thus check boundary values in three cases.

\medskip\noindent\emph{Case 1. $r\ge4$ and $b\ge2$.}\ 
Here $A\le3<s_0$. At $(a,t,\varepsilon)=(rb+1,2,0)$,
$\Psi=b-1-\max\{0,(3-r)b^2+2b+1\}\ge0$, covering all mixed forests and
all-odd forests with unequal odd components. When all components are
equal odd paths, we have $a=rb$ and $t\ge3$, giving
$\Psi=2b-2-(r+1-2b)/(br-1)-A>2b-2-1/b-A$; this is nonnegative since
$A\le1$ for $b\ge3$ and $A=0$ for $b=2,r\ge5$.

The remaining case $H=4P_5$ ($r=4,b=2,a=8,\varepsilon=1$) is handled by
including $I$ from Lemma~\ref{lem:residual}: before $I$, the gap is
$2t-50/7$, while $I=3(21-\binom{t+1}{2})/13$ for $3\le t\le6$ and
$I=3(t-6)/2$ for $t\ge6$. The sum is concave on $[3,6]$ with endpoint
values $211/91$ and $34/7$, and increases for $t\ge6$.

\medskip\noindent\emph{Case 2. $r\ge4$ and $b=1$.}\ 
Since $H$ is not a union of copies of $P_3$, $a\ge r+1$. If $a\ge r+2$,
monotonicity gives $\Psi\ge2-\max\{0,5-r\}>0$. If $a=r+1$, then
$H=P_2\cup rP_3$ ($\varepsilon=0$) or $H=P_5\cup(r-1)P_3$
($\varepsilon=1$), so Lemma~\ref{lem:residual} applies with
$\xi=\varepsilon/2$. Putting $z=2r-3$ and $A=\max\{0,6-r\}$, the gap
with $I$ on $2+\varepsilon\le t\le z$ is
\[
 t-2+\frac{2\varepsilon}{r}-A
 +\frac{2-\varepsilon}{2r-\varepsilon}
 \left[\binom{2r-2}{2}-\binom{t+1}{2}\right],
\]
which is concave in $t$, while for $t\ge z$ it is increasing. For
$\varepsilon=0$, both endpoints equal $2r-5-A\ge1$; for $\varepsilon=1$,
they are $r-1+2/r-5/(2r-1)-A\ge11/14$ and $2r-5+2/r-A\ge3/2$, both
increasing for $r\ge4$. Meanwhile $s_0-A=t+1-A\ge1$.

\medskip\noindent\emph{Case 3. $r=3$.}\ 
Set $x=a-3b$. If $x\ge3$, Theorem~\ref{thm:rooted-deletion} at the
boundary $(a,t,\varepsilon)=(3b+3,2,0)$ has $A=1<s_0$ and
$\Psi=b+4-12b/(3b+2)>0$.

Suppose $x\le2$, and let $2(b+u)+1$ be the order of the second shortest
odd component, so $u\in\{0,1\}$ (if $u=1$, then $a\ge3b+2$, forcing
$x=2$ and $\varepsilon=1$). Let $L=H-P_{2b+1}-P_{2(b+u)+1}$. Then
$|L|=2(a-2b-u)+1$ and $|R|-|L|=b-x+t+2u-1\ge0$, so
Lemma~\ref{lem:dirac-path} embeds $L$ into $R$. Delete a $P_{2b+1}$ from
$G-R$ using Lemma~\ref{lem:ordered-deletion}(iv) with $w(y)=d_R(y)$ and
$D=a-1$ (introducing no error term $A$). The remaining graph outside $R$
must be $P_{2(b+u)+1}$-free, as a copy of $P_{2(b+u)+1}$ there would
combine with the deleted $P_{2b+1}$ and $L\subseteq R$ to form $H$.
Applying Lemma~\ref{lem:residual} with $\eta=b+u-1/2$ gives gap
\[
 bt+\left(\frac{5}{4}-\frac{u}{2}\right)a-4b-3u+\frac{7}{2}
 -\left(b+\frac{u}{2}-\frac{1}{4}\right)f.
\]
Using $f\le4/b$, this is at least $7b/4-1/2+1/b>0$ when $u=0$; when
$u=1$, using $t\ge3$, $a=3b+2$, and $f=10/(3b+1)$ yields
$[15(b-1)^2+19(b-1)+2]/[4(3b+1)]>0$, while $s_0\ge0$ covers the join
term. Every case contradicts $e(G)>\max\{K,J_a(n)+\varepsilon\}$.
\end{proof}

\begin{lemma}\label{lem:uniform-lines}
Let $r\ge3$, $b\ge2$, $a\ge rb$, $\varepsilon\in\{0,1\}$, and
$2\le d\le2b-1$. Put
\[
 m=2a+r,\quad K=\binom{m-1}{2},\quad
 f=\frac{r(r+1)-2\varepsilon}{a-1},\quad
 n_0=\frac{5a+4r-2+f}{2},
\]
so that $K=J_a(n_0)+\varepsilon$, and let
$M_0(x)=\max\{K,J_a(x)+\varepsilon\}$. For $1\le h\le d-1$, define
\[
 \sigma_h=d-1+\frac{h-1}{2},\qquad
 F_h(x)=\binom{m-d}{2}+(d-1)(a-1)+\sigma_h(x-m+1).
\]
Setting $h_+=\min\{b,d-1\}$ and $h_-=\min\{b-1,d-1\}$, we have for every
real $x$,
\begin{equation}\label{eq:uniform-lines}
 F_{h_+}(x)<M_0(x)+4,\qquad F_{h_-}(x)<M_0(x).
\end{equation}
\end{lemma}
\begin{proof}
Put $g_h=K-F_h(n_0)$. Direct subtraction gives
\begin{equation}\label{eq:gh}
 4g_h=a(2d-h-1)-2d(d-1)-2r(h-1)-(2d+h-3)f.
\end{equation}
Since $\partial(4g_h)/\partial a=2d-h-1+(2d+h-3)[r(r+1)-2\varepsilon]/(a-1)^2>0$
and setting $\varepsilon=0$ maximizes $f$, the minimum of $g_h$ occurs
on the boundary $a=rb,\varepsilon=0$. Along this boundary,
$f=r(r+1)/(br-1)$ satisfies $df/dr=(br^2-2r-1)/(br-1)^2<1/b$, whence
$d(4g_h)/dr>b(2d-h-1)-2(h-1)-(2d+h-3)/b$. This lower bound decreases in
$h$, so it is minimized at $h=h_+$, where it is piecewise linear in $d$
on $[2,b+1]$ and $[b+1,2b-1]$ with positive endpoint values $2b-2/b$ and
$(b-1)(b^2-1)/b$. Thus $g_{h_+}$ and $g_{h_-}$ both increase with $r$.

Setting $r=3$, $u=b-2\ge0$, and $D=4(3b-1)$, the expression
\eqref{eq:gh} is a concave quadratic in $d$ on each subinterval where
$h_+$ or $h_-$ is constant or equals $d-1$. Evaluating at the endpoints
gives
\[
\begin{aligned}
 Dg_1\big|_{d=2} &= 2(3u+1)(3u+8),\\
 4\bigl(g_b\big|_{d=b+1}+4\bigr) &= u^2-u+4,\\
 D\bigl(g_b\big|_{d=2b-1}+4\bigr) &= 3u^3+8(u-1)^2+9u+12,\\
 Dg_{b-1}\big|_{d=b} &= 3u(u-1)(u+1)+5(u-1)^2+u+11,\\
 Dg_{b-1}\big|_{d=2b-1} &= (u+6)(3u^2-u+2),
\end{aligned}
\]
all strictly positive for $u\ge0$. By concavity, $g_{h_+}>-4$ and
$g_{h_-}>0$ for all $2\le d\le2b-1$. Finally,
$0<\sigma_h\le(5b-5)/2<a-1$, so $M_0(x)-F_h(x)$ attains its minimum at
$x=n_0$, yielding \eqref{eq:uniform-lines}.
\end{proof}

\begin{proposition}\label{prop:exists-low-core-path}
Let $H$ be a linear forest with $r\ge3$ odd components, no isolated
vertices, and shortest odd component $P_{2b+1}$, and suppose $H$ is not
a union of copies of $P_3$. Put $a=\nu(H)$, $m=2a+r$, and
$K=\binom{m-1}{2}$, and let $\varepsilon=1$ if all components of $H$ are
odd and $\varepsilon=0$ otherwise. Define
$M(n)=\max\{K,J_a(n)+\varepsilon\}$, except that $K$ is replaced by
$B_m(n)$ when $\varepsilon=1$, $r\in\{3,4\}$, and $b\ge2$. If $G$ is a
connected $H$-free graph of order $n\ge m$ with $e(G)>M(n)$, then $G$
contains a copy of $P_{2b+1}$ using at most $b$ vertices of the $a$-core
$R$ from Proposition~\ref{prop:core}.
\end{proposition}
\begin{proof}
Suppose for a contradiction that every $P_{2b+1}$ in $G$ uses at least
$b+1$ vertices of $R$. Let $Y=V(G)\setminus R$ with its deletion order
from Proposition~\ref{prop:core}, write $c=|R|=a+t=m-d$ and $N=|Y|$, and
let $\lambda,\Delta$ be as before.

If $b=1$, then $a\ge r+1$. Because every $P_3$ uses at least two
vertices of $R$, $Y$ is independent and every vertex of $R$ has at most
one neighbor in $Y$, giving $e(G)\le\binom{c+1}{2}-\lambda$. From
$e(G)>K$ we obtain $c=m-1$ and $\lambda<c$. Pick a vertex of $Y$
adjacent to $R$ and add it to $R$; the resulting $m$-vertex graph has no
Hamiltonian path (else it contains $H$). Proposition~\ref{prop:posa-structure}
with $d=1$ then yields
$\lambda\ge\min\{(a-1)(a+2r)/2,(c^2-1)/8\}>c$ (using $a\ge r+1$,
$r\ge3$, and $c\ge10$), a contradiction. Thus $b\ge2$, so $a\ge rb$ (and
$a\ge rb+1$ if $\varepsilon=0$).

Set $f=[r(r+1)-2\varepsilon]/(a-1)$, $n_0=(5a+4r-2+f)/2$,
$M_0(x)=\max\{K,J_a(x)+\varepsilon\}$, and
$T_d=\binom{m-d+1}{2}+(b-1)n_0-K$. By Lemma~\ref{lem:rooted-joint},
$e(G)\le\binom{c+1}{2}-\lambda+(b-1)n$. Since the slope $b-1$ lies
strictly between $0$ and $a-1$, comparing at $n_0$ forces
$\lambda<T_d$. Notice that $T_{d+1}-T_d=-(m-d)<0$, and at $d=2b$ we have
$T_{2b}=-(b-1)(3a-4b-f)/2<0$ because $f<(r+2)/b<a$. Hence
$1\le d\le2b-1<a$, and $N=n-c\ge m-(m-d)=d$.

\medskip\noindent\emph{Claim 1. If $r\ge5$, then for any $1\le h\le d$,
adjoining $h$ vertices of $Y$ to $R$ so that each has degree at least
$h$ in the induced subgraph yields a graph with a Hamiltonian path.}\ 
If not, Proposition~\ref{prop:posa-structure} (with $r-d+h$ in place of
$r$) and \eqref{eq:posa-endpoints} give
$\lambda\ge\min\{L_{d,h},D_{d,h}\}$, where
$L_{d,h}=(a-h)(2c+h-3a+1)/2$ and $D_{d,h}=(c-h)(c-h+2)/8$ (if the
parameter range of Proposition~\ref{prop:posa-structure} is empty, a
Hamiltonian path exists immediately). Writing $w=d-h\ge0$, the
differences $L_{d,h}-T_d$ and $D_{d,h}-T_d$ exceed their values at
$(d,h)=(1,1)$ by
\[
 a(h-1)+\frac{w(2a+2r+1-w)}{2}
 \quad\text{and}\quad
 \frac{(h-1)(m-w)}{2}+\frac{w(6m-2-3w)}{8},
\]
respectively, both nonnegative. At $(1,1)$, we have $D_{1,1}=m(m-2)/8$,
$8(L_{1,1}-D_{1,1})=r(4a-r-6)>0$, and
$8(D_{1,1}-T_1)=4a(a-5b)+4r(a-4b)+r^2+6r+8b-4(b-1)f$. Since $a\ge5b$ and
$(b-1)f<r+2$ (as $(r+2)(rb-1)-(b-1)r(r+1)=rb+r^2-2>0$), this gives
$8(D_{1,1}-T_1)>4a(a-5b)+4r(a-4b)+r^2+2r+8b-8>0$, contradicting
$\lambda<T_d$. In particular, taking $h=1$ with a pendant leaf at any
vertex $x\in R$ shows that $R$ has a Hamiltonian path starting at $x$.

\medskip\noindent\emph{Claim 2. One can choose $d$ vertices of $Y$ in
sequence so that each has at least $d$ neighbors among $R$ and the
earlier chosen vertices.}\ 
Observe first that when $d\ge2$, no $P_d$ in $G[Y]$ has an endpoint
adjacent to $R$: if $r\ge5$, it would extend via Claim~1 to $P_m$; if
$r\in\{3,4\}$ and $d\ge r$, then $|R|\le2a$ and Dirac's theorem gives a
Hamiltonian cycle in $R$, again yielding $P_m$; and if $r\in\{3,4\}$ and
$2\le d<r$, then $s:=r-d\in\{1,2\}$ and
$\lambda\le\binom{a+s+1}{2}-\varepsilon-1$, so
Lemma~\ref{lem:rooted-completion} gives a copy of $H$. (When $d=1$,
connectedness directly supplies a vertex of $Y$ adjacent to $R$.)

Select vertices greedily: set $C_0=\varnothing$, and at step $j$ pick
any $y\in Y\setminus C_{j-1}$ with $d_{R\cup C_{j-1}}(y)\ge d$, setting
$C_j=C_{j-1}\cup\{y\}$. Suppose this stops after $v<d$ steps with set
$C$, and put $Y'=Y\setminus C$. Note that $d\ge2$, and every vertex of
$C$ has a neighbor in $R$ (the $j$th chosen vertex has at least $d-j+1$
neighbors in $R$). Mark each vertex of $Y'$ that has a neighbor in
$R\cup C$; every component of $G[Y']$ contains a marked vertex.

No marked vertex can be an endpoint of a $P_{b+1}$ or a $P_d$ in
$G[Y']$. Indeed, for a path $y_1\cdots y_{b+1}$ in $G[Y']$ with $y_1$
marked, if $y_1$ has a neighbor $p\in R$ we prepend a $b$-vertex path in
$R$ ending at $p$, while if $y_1$ has a neighbor $w\in C$ we pick
$x\in N_R(w)$ and prepend a $(b-1)$-vertex path in $R$ ending at $x$
followed by $w$ (both exist since $\delta(R)\ge a\ge b$); either way we
get a $P_{2b+1}$ using at most $b$ vertices of $R$. Similarly, a path
$y_1\cdots y_d$ in $G[Y']$ with $y_1$ marked gives a $P_d$ in $G[Y]$
with an endpoint adjacent to $R$ (either $y_1\cdots y_d$ directly, or
$wy_1\cdots y_{d-1}$ for $w\in N_C(y_1)$), contradicting our initial
observation.

Put $h=\min\{b,d-1\}$ and let $\sigma_h,F_h$ be as in
Lemma~\ref{lem:uniform-lines}, with $z$ denoting the number of unmarked
vertices in $Y'$. By Lemma~\ref{lem:marked-root} and the stopping rule,
$e(R\cup C,Y')+e(Y')\le\sigma_h(|Y|-v)-(d-3/2)z$. Ordering $C$ by the
core-deletion order from Proposition~\ref{prop:core} gives
$e(R,C)+e(C)\le v(a-1)$, so
$e(G)\le F_h(n)-\lambda-(d-1-v)(a-1-\sigma_h)$. By the first inequality
in \eqref{eq:uniform-lines} and integrality, $\lambda\ge3$ contradicts
$e(G)>M_0(n)$, and $v=0$ likewise fails because
$(d-1)(a-1-\sigma_h)\ge4$ (it is $a-2\ge4$ for $d=2$, and
$\sigma_h\le3b-3\le a-3$ for $d\ge3$). Hence $v\ge1$ and $\lambda\le2$.

We can now tighten to $h=\min\{b-1,d-1\}$. When $d\ge b+1$, no marked
vertex can be an endpoint of a $P_b$ in $G[Y']$: write such a path as
$y_1\cdots y_b$. If $y_1$ has a neighbor $w\in C$, prepending a
$b$-vertex path in $R$ ending at a neighbor of $w$ and then $w$ gives a
$P_{2b+1}$ with $b$ vertices in $R$. If $y_1$ has a neighbor $p\in R$,
let $w$ be the first vertex chosen in $C$ (so $d_R(w)\ge d\ge b+1$).
When $pw\in E(G)$, pick $q\in N_R(w)\setminus\{p\}$ and a $(b-1)$-vertex
path in $R-p$ ending at $q$, giving $r_1\cdots r_{b-1}wpy_1\cdots y_b$.
When $pw\notin E(G)$, build a $(b-1)$-vertex path $p=r_1\cdots r_{b-1}$
in $R$; at least three neighbors of $w$ in $R$ lie outside this path,
and since $\lambda\le2$, $r_{b-1}$ is adjacent to one of them, say $q$,
yielding $y_b\cdots y_1r_1\cdots r_{b-1}qw$. Both give a forbidden
$P_{2b+1}$. Applying the second inequality in \eqref{eq:uniform-lines}
now yields $e(G)<M_0(n)$, a contradiction. This proves Claim~2.

If $r\ge5$, then $R$ together with the $d$ vertices from Claim~2 has a
Hamiltonian path by Claim~1, contradicting that $G$ is $H$-free. Assume
therefore that $r\in\{3,4\}$.

\medskip\noindent\emph{Claim 3. If $r=3$, then $d\le b$ (with $d=1$ when
$b=2$); if $r=4$, then $b\ge3$ and $d\le\lfloor b/2\rfloor$.}\ 
Let $U$ be the $d$-set from Claim~2. Since $Q=G[R\cup U]$ has no
Hamiltonian path, \eqref{eq:budget23}--\eqref{eq:budget41} imply
$\lambda\ge D_d:=\binom{a+3-d}{2}+r-4$. Set $\Delta_d=D_d-T_d$; since
$\lambda<T_d$, we must have $\Delta_d<0$. Computing directly,
\[
 \Delta_d=\binom{a+3-d}{2}+r-4-\binom{2a+r-d+1}{2}
 -(b-1)\frac{5a+4r-2+f}{2}+\binom{2a+r-1}{2},
\]
which satisfies
\begin{align}
 \Delta_{d+1}-\Delta_d &= a+r-2, \label{eq:finite-diff}\\
 \frac{\partial\Delta_d}{\partial a}
 &= a-\frac{5b}{2}+d+1+\frac{(b-1)[r(r+1)-2\varepsilon]}{2(a-1)^2}>0,
 \notag\\
 \Delta_d\big|_{\varepsilon=1}-\Delta_d\big|_{\varepsilon=0}
 &= \frac{b-1}{a-1}\ge0. \notag
\end{align}
Thus $\Delta_d$ increases in $d$, $a$, and $\varepsilon$. Checking the
minimum boundaries gives
\begin{align*}
 \Delta_{b+1}\big|_{r=3,a=3b,\varepsilon=0} &= \frac{6b^2+b+3}{3b-1}>0,
 & \Delta_b\big|_{r=3,a=3b+1,\varepsilon=0} &= \frac{b^2+3b+4}{2b}>0,\\
 \Delta_{(b+1)/2}\big|_{r=4,a=4b,\varepsilon=0} &= \frac{6(b+1)}{4b-1}>0,
 & \Delta_1\big|_{r=4,a=9,b=2,\varepsilon=0} &= \frac{13}{4}>0,
\end{align*}
which by \eqref{eq:finite-diff} yields $d\le b$ for $r=3$ ($d\le b-1$
when $a\ge3b+1$), $d\le\lfloor b/2\rfloor$ for $r=4$ (since every
integer $d\ge\lfloor b/2\rfloor+1$ is at least $(b+1)/2$), and excludes
$r=4,b=2,a\ge9$.

Two all-odd boundary points remain. For $(r,a,b,d)=(3,6,2,2)$,
Lemma~\ref{lem:rooted-joint} and $\lambda\ge D_d$ give $e(G)\le n+71$,
which exceeds $M_0(n)$ only at $n=21$, where $B_{15}(21)=94>92$. For
$(r,a,b,d)=(4,8,2,1)$, they give $e(G)\le n+145$, exceeding $M_0(n)$
only at $n\in\{27,28\}$, where $B_{20}(n)=175>n+145$. Both are ruled
out, proving Claim~3.

\medskip\noindent\emph{Claim 4. One can choose $U$ of size $d$ in
sequence so that each vertex has at least
$\kappa=\lceil(a+d+2r-5)/2\rceil$ neighbors among $R$ and the earlier
chosen vertices.}\ 
Retaining $\lambda\ge D_d$, run the greedy selection with threshold
$\kappa$ instead of $d$, and suppose it stops after $v<d$ vertices
$w_1,\ldots,w_v$, forming $C$. Each $w_j$ has at least
$\kappa-j+1\ge\kappa-d+1\ge1$ neighbors in $R$. Let $Y'=Y\setminus C$,
$q=|Y'|$, and mark every vertex of $Y'$ adjacent to $R\cup C$; let $z$
be the number of unmarked vertices. As in Claim~2, every component of
$G[Y']$ has a marked vertex and no marked vertex is an endpoint of a
$P_{b+1}$ in $G[Y']$. By Lemma~\ref{lem:marked-root},
$e(Y')\le(b-1)q/2+z/2$, while the stopping rule gives at most $\kappa-1$
edges from each marked vertex to $R\cup C$ and none from unmarked
vertices. Putting $\sigma=\kappa+(b-3)/2$, we have
$e(R\cup C,Y')+e(Y')\le\sigma q-(\kappa-3/2)z\le\sigma q$. Combining
this with $e(R,C)+e(C)\le v(a-1)$, $\lambda\ge D_d$, and $v\le d-1$
yields
\begin{equation}\label{eq:stopping-line}
 e(G)\le S_d(n):=\binom{m-d}{2}-D_d+(d-1)(a-1-\sigma)+\sigma(n-m+d).
\end{equation}

By Claim~3, $0<\sigma<a-1$, so $M_0(n)-S_d(n)$ is minimized at $n_0$.
Put $g=K-S_d(n_0)$. Because
$\partial g/\partial\sigma=-(a^2+2ar-a-2\varepsilon+r^2-r)/[2(a-1)]<0$
and $\sigma\le(a+b+d+2r-7)/2$, substituting into \eqref{eq:stopping-line}
gives
\[
 4g\ge
 \begin{cases}
  Q_3+\dfrac{2\varepsilon(a+b+d-1)}{a-1}, & r=3,\\[6pt]
  Q_4+\dfrac{2\varepsilon(a+b+d+1)}{a-1}, & r=4,
 \end{cases}
\]
where
\[
\begin{aligned}
 Q_3 &= a^2-a(b+d-1)-6b-2d-10-\frac{12(b+d)}{a-1},\\
 Q_4 &= a^2-a(b+d+3)-8b-32-\frac{20(b+d+2)}{a-1}.
\end{aligned}
\]
Both $Q_3$ and $Q_4$ decrease in $d$ and increase in $a$:
\begin{align*}
 \frac{\partial Q_3}{\partial d} &= -a-2-\frac{12}{a-1}<0, &
 \frac{\partial Q_3}{\partial a} &= 2a-b-d+1+\frac{12(b+d)}{(a-1)^2}>0,\\
 \frac{\partial Q_4}{\partial d} &= -a-\frac{20}{a-1}<0, &
 \frac{\partial Q_4}{\partial a} &= 2a-b-d-3+\frac{20(b+d+2)}{(a-1)^2}>0,
\end{align*}
so it suffices to check maximum $d$ and minimum $a$ allowed by Claim~3.

For $r=3$: when $a\ge3b+1$ and $d\le b-1$, we get
$g\ge(3b-1)(b^2+b-4)/(4b)>0$; when $a=3b$ and $d\le b$, we get
$g\ge(3b+5)(3b^2-11b+2)/(4(3b-1))>0$ for $b\ge4$; when $(a,b)=(9,3)$
and $d\le2$, $g\ge11/8$; and when $(a,b,d)=(6,2,1)$, exact values
$\kappa=4$, $\sigma=7/2$, $n_0=21$ give $g=5/2$. The lone remaining
point is $(a,b,d,\varepsilon)=(9,3,3,1)$, where $\kappa=\sigma=7$ and
$g=-7/8$ at $n_0=225/8$; since $M_0(n)-S_d(n)$ equals $g+7(n_0-n)$ for
$n\le28$ and $g+(n-n_0)$ for $n\ge29$, it vanishes at $n=28,29$ and is
positive at every other integer, so $S_d(n)\le M_0(n)$ for all integer
$n$.

For $r=4$: when $b\ge4$, at $a=4b,d=b/2$ we have
$g\ge(b+1)(20b^2-65b-4)/(2(4b-1))>0$; when $b=3$, we have $d=1$ and
$g\ge3$ at $a=13$, while at $a=12$ exact values $\kappa=8,f=18/11$ give
$g=49/11>0$. In every case $e(G)\le M_0(n)$, contradicting $e(G)>M(n)$
and proving Claim~4.

Finally, let $Q=G[R\cup U]$ for the set $U$ from Claim~4. Since $Q$
cannot have a Hamiltonian path, and $t=a+r-d$ satisfies $t\ge8$ for
$r=3$ and $t\ge15$ for $r=4$ by Claim~3,
Proposition~\ref{prop:posa-structure} provides a subset $W\subseteq V(Q)$
of size $a+3$ with $e(Q[W])\le1$. Applying
Theorem~\ref{thm:factor-from-W} then produces a spanning copy of $H$ in
$Q$, contradicting that $G$ is $H$-free.
\end{proof}

\subsection{Proof of the connected case}

\begin{theorem}\label{thm:connected}
Let $H$ be a nonempty linear forest whose components all have order at
least two, and assume $H$ is not a union of copies of $P_3$. Put
$a=\nu(H)$ and $m=|H|=2a+O$, where $O$ is the number of odd components,
and let $\varepsilon=1$ if all components of $H$ are odd and
$\varepsilon=0$ otherwise. If $G$ is a connected $H$-free graph of order
$n\ge m$, then $e(G)\le\max\{B_m(n),J_a(n)+\varepsilon\}$. Moreover, if
$H$ has a $P_2$ component or if $O\ge5$, then $B_m(n)$ may be replaced
by $\binom{m-1}{2}$.
\end{theorem}
\begin{proof}
We proceed by induction on $O$. The base range $O\le2$ is settled by
Theorem~\ref{thm:at-most-two-odd}.

Now let $O\ge3$, pick a shortest odd component $P_{2b+1}$ of $H$, and
set $H'=H-P_{2b+1}$. Note that $H'$ is not a union of copies of $P_3$:
if $b=1$, then $H$ contains a component other than $P_3$, which remains
in $H'$; if $b\ge2$, then $H'$ still contains at least two odd
components of order at least $5$. By the induction hypothesis, the
connected bound holds for $H'$ at all orders at least $|H'|$, with the
stronger constant term $\binom{|H'|-1}{2}$ whenever $H'$ has a $P_2$
component.

Suppose for a contradiction that $G$ violates the theorem, so in
particular $e(G)>\max\{\binom{m-1}{2},J_a(n)+\varepsilon\}$. By
Proposition~\ref{prop:core}, $G$ has a connected $a$-core $R$ satisfying
all properties listed there. Proposition~\ref{prop:no-low-core-path}
then implies that no copy of $P_{2b+1}$ in $G$ uses at most $b$ vertices
of $R$, which contradicts Proposition~\ref{prop:exists-low-core-path}.
\end{proof}

\section{Proof of Theorem~\ref{thm:main}}\label{sec:main}

For a nonempty linear forest $L$ without isolated vertices, let
$\ell(L)$ be the order of its shortest component, and set
$\varepsilon(L)=1$ if every component of $L$ is odd and
$\varepsilon(L)=0$ otherwise. Whenever $L$ is neither a matching nor a
union of copies of $P_3$, Theorem~\ref{thm:connected} gives the
connected bound
\begin{equation}\label{eq:connected-bound}
 e(X)\le\max\bigl\{\Pi(x;|L|,\ell(L)),\,J_{\nu(L)}(x)+\varepsilon(L)\bigr\}
 \qquad(x=|X|\ge|L|)
\end{equation}
for every connected $L$-free graph $X$. Indeed, for $\ell(L)=2$ the
first term of Theorem~\ref{thm:connected} is $\binom{|L|-1}{2}=\Pi(x;|L|,2)$,
for $\ell(L)=3$ it is $B_{|L|}(x)=\Pi(x;|L|,3)$, and for $\ell(L)\ge4$
we verify below that $B_{|L|}(x)\le\Pi(x;|L|,\ell(L))$. For matchings,
\eqref{eq:connected-bound} also holds by the Erd\H{o}s--Gallai matching
theorem~\cite{EG}, while unions of copies of $P_3$ are governed by
\eqref{eq:triples} and single paths satisfy
$\ex(z,P_k)=\Pi(z;k,k)$ (including $k=2$).

\subsection{The disconnected case}

Write $\mathcal{C}(x)=x(x-1)/2$ for any real $x$. Fix an integer
$k=2h+\eta\ge2$ with $\eta\in\{0,1\}$, and set $c=k-1$,
$\lambda=(c-1)/2$, $\mathcal{Q}(m)=m(m+c)/2$, $\rho=c^2/8$, and
$f_k(z)=\Pi(z;k,k)$. Writing $z=uc+v$ with $0\le v<c$, we have
$f_k(z)=u\mathcal{C}(c)+\mathcal{C}(v)=\lambda z-v(c-v)/2\le\lambda z$.
For $x=uc+v$ and $y=u'c+w$, direct subtraction gives
$f_k(x+y)-f_k(x)-f_k(y)=vw$ when $v+w<c$ and $(c-v)(c-w)$ when
$v+w\ge c$, so $f_k$ is superadditive (including $f_2\equiv0$).

When $k\ge4$, we have $\mathcal{C}(c)\ge c$ and
$\mathcal{C}(v)\ge v-1$, whence $f_k(z)\ge z-1\ge\lfloor z/2\rfloor$ for
all $z\ge1$ (and $f_k(0)=0$), confirming that $B_s(x)\le\Pi(x;s,k)$
whenever $x\ge s\ge k\ge4$. For $P_m(x):=\Pi(x;m+k,k)$ and $t\ge0$,
\begin{equation}\label{eq:prefix-id}
 P_m(m+t)=\mathcal{C}(m)+m\min\{t,c\}+f_k(t),
\end{equation}
which immediately yields
\begin{equation}\label{eq:affine-prefix}
 P_m(x)\le\lambda x+\mathcal{Q}(m)\quad(x\ge m),\qquad
 P_m(x)\ge\lambda x+\mathcal{Q}(m)-\rho\quad(x\ge m+c),
\end{equation}
with slack $(c-t)(m+t/2)$ in the first inequality for $0\le t\le c$.
Together with the superadditivity of $f_k$, \eqref{eq:prefix-id} also
implies $P_m(w)+f_k(z)\le P_m(w+z)$ for all $w\ge m$ and $z\ge0$.

Yuan and Zhang~\cite[Lemma~2.3]{YZ} used superadditivity and a
component-allocation argument for forests with at most one odd
component. The following theorem extends this reduction to arbitrary
linear forests, accommodating $P_2$ components and $tP_3$ subforests.

\begin{theorem}\label{thm:disconnected-reduction}
Let $F$ be a nonempty linear forest without isolated vertices. Assume
that \eqref{eq:connected-bound} holds for every subforest $L$ of $F$
formed by whole components, except possibly when $L$ is a matching or a
union of copies of $P_3$. Combined with the exact formulas for
matchings, unions of copies of $P_3$, and single paths, this implies
\eqref{eq:formula} whenever $F$ is not a union of copies of $P_3$, and
\eqref{eq:triples} when $F$ is a union of copies of $P_3$.
\end{theorem}
\begin{proof}
We assume $n\ge|F|$ (the lower bounds and $n<|F|$ were settled in
Section~\ref{sec:intro}) and induct on the number of components of $F$.
The theorem is already known when $F$ is a single path, a matching, or a
union of copies of $P_3$.

Let $P_k$ be a shortest component of $F$, and write $F=F^-\cup P_k$. Put
$m=|F|$, $a=\nu(F)$, and $h=\lfloor k/2\rfloor$.

Suppose first that $G$ is $F^-$-free. By induction, $e(G)$ satisfies the
formula for $F^-$. Every prefix term of $F^-$ is also a prefix term of
$F$. When $F^-$ is not a union of copies of $P_3$, its join term is
strictly smaller than that of $F$, since for $n\ge m\ge2a$,
\[
 J_a(n)+\varepsilon(F)-J_{a-h}(n)-\varepsilon(F^-)
 =h\left(n-a+\frac{h+1}{2}\right)+\varepsilon(F)-\varepsilon(F^-)>0.
\]
If $F^-=tP_3$, then because $F$ is not a union of copies of $P_3$ and
$P_k$ is a shortest component, we must have $k=2$. The first branch
$B_{3t}(n)$ of \eqref{eq:triples} is the prefix term for the last $P_3$
component of $F$, while the second branch is at most $J_{t+1}(n)$,
because
\[
 J_{t+1}(n)-\left[(t-1)n-\mathcal{C}(t)
 +\left\lfloor\frac{n-t+1}{2}\right\rfloor\right]
 =n-t-\left\lfloor\frac{n-t+1}{2}\right\rfloor\ge0
\]
(including $t=1$). Thus $e(G)$ satisfies \eqref{eq:formula} whenever $G$
is $F^-$-free.

Now fix a copy of $F^-$ in $G$. Let $X_1,\ldots,X_s$ be the connected
components of $G$ meeting this copy, and let $W_i$ be the union of those
components of $F^-$ embedded in $X_i$. Then each $X_i$ is
$(W_i\cup P_k)$-free (otherwise replacing the embedding in $X_i$ would
produce $F$), and the union $U$ of all remaining components of $G$ is
$P_k$-free. Set $a_i=\nu(W_i)$, $m_i=|W_i|$, and $x_i=|X_i|$. Since
every component of $W_i$ has order at least $k$, we have $a_i\ge h$,
$2a_i\le m_i\le3a_i$, and $x_i\ge m_i$.

If $s=1$, then $W_1\cup P_k=F$ is neither a matching nor a union of
copies of $P_3$. Using \eqref{eq:connected-bound} when $x_1\ge m$ and
the complete-graph bound when $x_1<m$, we obtain
$e(X_1)\le\max\{P_{m_1}(x_1),J_a(x_1)+\varepsilon(F)\}$. Adding
$e(U)\le f_k(|U|)$ is absorbed by $P_{m_1}$ via \eqref{eq:prefix-id} and
by $J_a+\varepsilon(F)$ because its slope $a-1$ is at least $\lambda$.

Assume henceforth that $s\ge2$. We claim that for every $1\le i\le s$,
\begin{equation}\label{eq:local-affine}
 e(X_i)\le\max\bigl\{\lambda x_i+\mathcal{Q}(m_i),\,
 (a_i+\lambda)x_i-\mathcal{C}(a_i+h)+1\bigr\}.
\end{equation}
If $W_i\cup P_k$ is not a union of copies of $P_3$, this follows from
the complete-graph bound and \eqref{eq:affine-prefix} when
$x_i<m_i+k$, and from \eqref{eq:connected-bound} (or the matching bound)
when $x_i\ge m_i+k$, since $a_i+h-1\le a_i+\lambda$ and
$\varepsilon(W_i\cup P_k)\le1$. If $W_i\cup P_k=(t+1)P_3$ (so $k=3$,
$a_i=t$, $m_i=3t$), write $R(y)=y-2\lfloor y/2\rfloor\in\{0,1\}$; the
two branches of \eqref{eq:triples} for $(t+1)P_3$ rewrite as
\begin{align*}
 \mathcal{C}(3t+2)+\left\lfloor\frac{x-3t-2}{2}\right\rfloor
 &=\frac{1}{2}x+\mathcal{Q}(3t)-\frac{1}{2}R(x-3t),\\
 tx-\mathcal{C}(t+1)+\left\lfloor\frac{x-t}{2}\right\rfloor
 &=\left(t+\frac{1}{2}\right)x-\mathcal{C}(t+1)-\frac{t}{2}
  -\frac{1}{2}R(x-t),
\end{align*}
both at most the right-hand side of \eqref{eq:local-affine} (and for
$x\in\{3t,3t+1,3t+2\}$ the complete-graph bound suffices).

Put $A=\sum_i a_i=a-h$, $M=\sum_i m_i=m-k$, and $z=|U|$. When $k$ is
odd, we have $A\ge2h\ge4$ if $k\ge5$, and $A\ge3$ if $k=3$ (as some
$W_i$ contains a path longer than $P_3$ and at least one other $W_i$ is
nonempty). Also $n\ge m=M+c+1$. We claim that
\begin{align}
 &\sum_{i=1}^s\max\bigl\{\lambda x_i+\mathcal{Q}(m_i),\,
  (a_i+\lambda)x_i-\mathcal{C}(a_i+h)+1\bigr\}+\lambda z
  \notag\\
 &\hspace{35mm}<\max\bigl\{P_M(n),\,J_{A+h}(n)\bigr\}.
 \label{eq:sum-claim}
\end{align}
To prove \eqref{eq:sum-claim}, put $S=\sum_i\mathcal{Q}(m_i)$,
$E=n-M\ge c+1>0$, and
$H_i=\mathcal{Q}(m_i)+\mathcal{C}(a_i+h)-1-a_im_i$. Since
$\mathcal{Q}(x)-a_ix$ increases for $x\ge2a_i$, we have
$H_i\ge a_ic+\mathcal{C}(a_i+h)-1>0$. Writing $e_i=x_i-m_i\ge0$ and
$\phi_i(e)=\max\{0,a_ie-H_i\}$, direct algebra gives
\begin{align}
 &\max\bigl\{\lambda x_i+\mathcal{Q}(m_i),\,
  (a_i+\lambda)x_i-\mathcal{C}(a_i+h)+1\bigr\} \notag\\
 &\hspace{28mm}=\lambda x_i+\mathcal{Q}(m_i)+\phi_i(e_i).
 \label{eq:phi-identity}
\end{align}
Because $\sum_i e_i=E-z\le E$ and each $\phi_i$ is convex with
$\phi_i(0)=0$, we have $\phi_i(e_i)\le(e_i/E)\phi_i(E)$ and hence
$\sum_i\phi_i(e_i)\le\max_i\phi_i(E)$. The left-hand side of
\eqref{eq:sum-claim} is therefore at most
$\lambda n+S+\max_i\phi_i(E)$.

If $\max_i\phi_i(E)=0$, then \eqref{eq:affine-prefix} gives
$P_M(n)-\lambda n\ge\mathcal{Q}(M)-\rho$, whereas
$\mathcal{Q}(M)-S=\sum_{i<j}m_im_j\ge4h^2>\rho$, proving
\eqref{eq:sum-claim}.

Otherwise fix an index $i$ with $\phi_i(E)>0$, and write $\alpha=a_i$,
$p=A-\alpha$, $\mu=m_i$, $\beta=M-\mu$, $r_0=\alpha+p-\eta/2$,
$D=\mathcal{C}(\alpha+p+h)$, and $D_i=\mathcal{C}(\alpha+h)-1$. By the
superadditivity of $\mathcal{Q}$,
\[
 S+\alpha E-H_i=\alpha n+S-\mathcal{Q}(\mu)-\alpha\beta-D_i
 \le\alpha n+\mathcal{Q}(\beta)-\alpha\beta-D_i=:T(n).
\]
Set $G_P(n)=\mathcal{Q}(\mu+\beta)-\rho-T(n)$ (a lower bound for
$P_M(n)-\lambda n-T(n)$) and $G_J(n)=r_0n-D-T(n)=J_{A+h}(n)-\lambda n-T(n)$
(since $A+h-1-\lambda=r_0$). Taking the positive linear combination
with coefficients $r_0-\alpha=p-\eta/2>0$ and $\alpha>0$ eliminates $n$:
\[
 X:=(r_0-\alpha)G_P(n)+\alpha G_J(n)
 =(r_0-\alpha)[\mathcal{Q}(\mu+\beta)-\rho]-r_0\mathcal{Q}(\beta)
 +\alpha(r_0\beta-D)+r_0D_i.
\]
It suffices to show $X>0$. Over the rectangle $2\alpha\le\mu\le3\alpha$,
$2p\le\beta\le3p$ (with $\alpha,p\ge h$), we have
$\partial X/\partial\mu=(r_0-\alpha)(\mu+\beta+c/2)>0$ and
$\partial^2 X/\partial\beta^2=-\alpha<0$, so $X$ is minimized at
$\mu=2\alpha$ and $\beta\in\{2p,3p\}$. Denoting the value at $\beta=jp$
by $X_j$ ($j\in\{2,3\}$) and substituting $c=2h+\eta-1$, $\rho=c^2/8$,
$D$, and $D_i$, we obtain
\begin{align}
 8X_2\big|_{\eta=0} &= 28\alpha p(\alpha+p)-8\alpha-9p, \notag\\
 8X_2\big|_{\eta=1} &= 28\alpha p(\alpha+p)-10\alpha^2-12\alpha h
  -24\alpha p-6\alpha-4hp+2h-8p+4, \label{eq:X2-odd}\\
 8(X_3-X_2) &= 4\alpha p(2\alpha-2h+p+1-4\eta). \label{eq:endpoint-diff}
\end{align}
For $\eta=0$, \eqref{eq:endpoint-diff} is nonnegative (as $\alpha\ge h$)
and $8X_2\ge56\alpha p-8\alpha-9p\ge39$, so $X_2,X_3>0$. For $\eta=1$,
we have $\alpha+p=A\ge3$ and $8(X_3-X_2)\ge-8\alpha p$. In
\eqref{eq:X2-odd}, bounding $10\alpha^2\le10\alpha^2p$,
$12\alpha h\le12\alpha p$, $4hp\le4\alpha p$, $6\alpha+8p\le14\alpha p$,
and $2h+4\ge6$ (and subtracting $8\alpha p$ to cover $j=3$ as well)
gives $8X_j\ge\alpha p(18\alpha+28p-62)+6>0$ for $j\in\{2,3\}$, since
$\alpha,p\ge1$ and $\alpha+p\ge3$ imply $18\alpha+28p-62\ge2$. Thus
$X>0$, establishing \eqref{eq:sum-claim}.

Finally, combining \eqref{eq:local-affine}, $e(U)\le f_k(z)\le\lambda z$,
and \eqref{eq:sum-claim} gives
$e(G)<\max\{\Pi(n;m,k),J_a(n)\}\le\max\{\Pi(n;m,k),J_a(n)+\varepsilon(F)\}$,
closing the induction.
\end{proof}

\begin{proof}[Completion of the proof of Theorem~\ref{thm:main}]
When $H=tP_3$, \eqref{eq:triples} is the theorem of Campos and
Lopes~\cite{CL} and Yuan and Zhang~\cite{YZtriples}. Otherwise,
Theorem~\ref{thm:connected} supplies \eqref{eq:connected-bound} for all
non-exceptional connected subforests, and
Theorem~\ref{thm:disconnected-reduction} then yields the upper bound in
\eqref{eq:formula} for all $n\ge m$. Matching lower-bound constructions
were verified in Section~\ref{sec:intro}, and for $0\le n<m$ the
complete graph $K_n$ is trivially $H$-free.
\end{proof}

\section*{Declaration of competing interest}
The authors declare that they have no known competing financial interests or personal relationships that could have appeared to influence the work reported in this paper.

\section*{Data availability}
No data were used for the research described in this article.

\section*{Declaration on the use of AI}
The authors used ChatGPT (OpenAI) to assist in discussing proof strategies, checking proofs, and improving exposition. The authors reviewed and verified the mathematical content and take full responsibility for the manuscript.

\small
\begin{thebibliography}{99}
\setlength{\itemsep}{-1pt}

\bibitem{BabuDiwan}
Ch.~S. Babu and A.~A. Diwan,
\emph{Degree conditions for forests in graphs},
Discrete Mathematics \textbf{301} (2005), nos.~2--3, 228--231.

\bibitem{BGLS}
P.~N. Balister, E. Gy\H{o}ri, J. Lehel and R.~H. Schelp,
\emph{Connected graphs without long paths},
Discrete Mathematics \textbf{308} (2008), no.~19, 4487--4494.

\bibitem{BielakKieliszek}
H. Bielak and S. Kieliszek,
\emph{The Tur\'an number of the graph $2P_5$},
Discussiones Mathematicae Graph Theory \textbf{36} (2016), 683--694.

\bibitem{BondyChvatal}
J.~A. Bondy and V. Chv\'atal,
\emph{A method in graph theory},
Discrete Mathematics \textbf{15} (1976), 111--135.

\bibitem{BushawKettle}
N. Bushaw and N. Kettle,
\emph{Tur\'an numbers of multiple paths and equibipartite forests},
Combinatorics, Probability and Computing \textbf{20} (2011), no.~6,
837--853.

\bibitem{CL}
V. Campos and R. Lopes,
\emph{A proof for a conjecture of Gorgol},
Discrete Applied Mathematics \textbf{245} (2018), 202--207.

\bibitem{ChenZhang}
M.-Z. Chen and X.-D. Zhang,
\emph{Erd\H{o}s--Gallai stability theorem for linear forests},
Discrete Mathematics \textbf{342} (2019), no.~3, 904--916.

\bibitem{CTW}
Y. Chen, F. Tian and B. Wei,
\emph{Degree sums and path-factors in graphs},
Graphs and Combinatorics \textbf{17} (2001), no.~1, 61--71.

\bibitem{CYYZ}
Y.-H. Chen, J.-B. Yang, L.-T. Yuan and P. Zhang,
\emph{Exact generalized Tur\'an numbers for even linear forests},
Discrete Mathematics \textbf{347} (2024), no.~7, Article~113974.

\bibitem{ChibaTsugaki}
S. Chiba and M. Tsugaki,
\emph{A degree sum condition for the existence of a path-factor},
Ars Combinatoria \textbf{113} (2014), 441--450.

\bibitem{DHZ}
J. Deng, J. Hou and Q. Zeng,
\emph{The Tur\'an number of three disjoint paths},
Discussiones Mathematicae Graph Theory \textbf{44} (2024), no.~4,
1513--1537.

\bibitem{Diestel}
R. Diestel,
\emph{Graph Theory}, 6th ed.,
Graduate Texts in Mathematics 173, Springer, 2025.

\bibitem{Dirac}
G.~A. Dirac,
\emph{Some theorems on abstract graphs},
Proceedings of the London Mathematical Society (3) \textbf{2} (1952),
69--81.

\bibitem{DNYZ}
M. Dong, B. Ning, L.-T. Yuan and X.-D. Zhang,
\emph{Exact Tur\'an numbers of two vertex-disjoint paths},
arXiv preprint arXiv:2511.01509v1, 2025.

\bibitem{EO}
Y. Egawa and K. Ota,
\emph{Vertex-disjoint paths in graphs},
Ars Combinatoria \textbf{61} (2001), 23--31.

\bibitem{EnomotoOta}
H. Enomoto and K. Ota,
\emph{Partitions of a graph into paths with prescribed endvertices and lengths},
Journal of Graph Theory \textbf{34} (2000), no.~2, 163--169.

\bibitem{EG}
P. Erd\H{o}s and T. Gallai,
\emph{On maximal paths and circuits of graphs},
Acta Mathematica Academiae Scientiarum Hungaricae \textbf{10} (1959),
337--356.

\bibitem{FangYou}
X. Fang and L. You,
\emph{The Tur\'an number of $P_9\cup P_7$},
Computational and Applied Mathematics \textbf{43} (2024), no.~5,
Article~291.

\bibitem{FS}
R.~J. Faudree and R.~H. Schelp,
\emph{Path Ramsey numbers in multicolorings},
Journal of Combinatorial Theory, Series B \textbf{19} (1975), no.~2,
150--160.

\bibitem{FengHu}
L. Feng and Y. Hu,
\emph{The Tur\'an number of the graph $3P_5$},
Filomat \textbf{34} (2020), no.~10, 3395--3410.

\bibitem{FKL}
Z. F\"uredi, A. Kostochka and R. Luo,
\emph{A stability version for a theorem of Erd\H{o}s on nonhamiltonian graphs},
Discrete Mathematics \textbf{340} (2017), no.~11, 2688--2690.

\bibitem{Johansson}
R. Johansson,
\emph{An El-Zah\'ar type condition ensuring path-factors},
Journal of Graph Theory \textbf{28} (1998), no.~1, 39--42.

\bibitem{Kawarabayashi}
K. Kawarabayashi,
\emph{Graph partition into paths containing specified vertices},
Discrete Mathematics \textbf{248} (2002), 271--277.

\bibitem{Kopylov}
G.~N. Kopylov,
\emph{On maximal paths and cycles in a graph},
Doklady Akademii Nauk SSSR \textbf{234} (1977), no.~1, 19--21.

\bibitem{LMY}
J. Lai, T. Ma and W. Yan,
\emph{The Tur\'an number of vertex disjoint paths of lengths one and two},
Discrete Mathematics \textbf{349} (2026), no.~2, Article~114764.

\bibitem{LanQinShi}
Y. Lan, Z. Qin and Y. Shi,
\emph{The Tur\'an number of $2P_7$},
Discussiones Mathematicae Graph Theory \textbf{39} (2019), no.~4,
805--814.

\bibitem{LLP}
B. Lidick\'y, H. Liu and C. Palmer,
\emph{On the Tur\'an number of forests},
Electronic Journal of Combinatorics \textbf{20} (2013), no.~2,
Paper~P62, 13~pp.

\bibitem{MO}
C. Magnant and K. Ozeki,
\emph{Partitioning graphs into paths or cycles of prescribed lengths},
Journal of Combinatorics \textbf{3} (2012), no.~2, 135--161.

\bibitem{Posa}
L. P\'osa,
\emph{A theorem concerning Hamilton lines},
Magyar Tudom\'anyos Akad\'emia Matematikai Kutat\'o Int\'ezet\'enek
K\"ozlem\'enyei \textbf{7} (1962), nos.~1--2, 225--226.

\bibitem{YZtriples}
L.-T. Yuan and X.-D. Zhang,
\emph{The Tur\'an number of disjoint copies of paths},
Discrete Mathematics \textbf{340} (2017), no.~2, 132--139.

\bibitem{YZ}
L.-T. Yuan and X.-D. Zhang,
\emph{Tur\'an numbers for disjoint paths},
Journal of Graph Theory \textbf{98} (2021), no.~3, 499--524.

\bibitem{ZZL}
H. Zhang, X. Zhao and M. Lu,
\emph{Tur\'an type problems for a fixed graph and a linear forest},
European Journal of Combinatorics \textbf{137} (2026), Article~104412.

\bibitem{ZhangWang}
L.-P. Zhang and L. Wang,
\emph{The Tur\'an numbers of special forests},
Graphs and Combinatorics \textbf{38} (2022), no.~3, Article~84.

\bibitem{ZhuChen}
X. Zhu and Y. Chen,
\emph{Generalized Tur\'an number for linear forests},
Discrete Mathematics \textbf{345} (2022), no.~10, Article~112997.

\end{thebibliography}
\end{document}
